\documentclass[11pt]{article}
\usepackage[T1]{fontenc}
\usepackage[utf8]{inputenc}
\usepackage{lmodern}
\usepackage[margin=1.04in]{geometry}
\usepackage{amsmath,amssymb,amsthm,mathtools}
\usepackage{aliascnt}
\usepackage{microtype}
\usepackage{mathrsfs}
\usepackage{enumitem}
\usepackage{booktabs,array}
\usepackage{xcolor}
\usepackage{hyperref}
\usepackage[capitalize,noabbrev]{cleveref}
\hypersetup{colorlinks=true,linkcolor=blue!35!black,citecolor=blue!35!black,urlcolor=blue!40!black,
 pdftitle={Rational-Exponential Iterated Integrals and Quantum Polylogarithms},
 pdfauthor={Christopher D. Long},
 pdfsubject={Absolute fixed-parameter functional completeness, quantum polylogarithms, difference equations, shuffle algebras, pentagon crossed products}}
\setlength{\emergencystretch}{3em}
\numberwithin{equation}{section}
\allowdisplaybreaks[2]
\newtheorem{theorem}{Theorem}[section]
\newaliascnt{proposition}{theorem}
\newtheorem{proposition}[proposition]{Proposition}
\aliascntresetthe{proposition}
\newaliascnt{lemma}{theorem}
\newtheorem{lemma}[lemma]{Lemma}
\aliascntresetthe{lemma}
\newaliascnt{corollary}{theorem}
\newtheorem{corollary}[corollary]{Corollary}
\aliascntresetthe{corollary}
\theoremstyle{definition}
\newaliascnt{definition}{theorem}
\newtheorem{definition}[definition]{Definition}
\aliascntresetthe{definition}
\theoremstyle{remark}
\newaliascnt{remark}{theorem}
\newtheorem{remark}[remark]{Remark}
\aliascntresetthe{remark}
\newaliascnt{example}{theorem}
\newtheorem{example}[example]{Example}
\aliascntresetthe{example}
\newcommand{\C}{\mathbb C}
\newcommand{\Z}{\mathbb Z}
\newcommand{\Q}{\mathbb Q}
\newcommand{\R}{\mathbb R}
\newcommand{\N}{\mathbb Z_{\geq0}}
\newcommand{\dd}{\,d}
\newcommand{\ii}{\mathrm i}
\newcommand{\sh}{\operatorname{sh}}
\newcommand{\Mer}{\operatorname{Mer}}
\newcommand{\Sh}{\operatorname{Sh}}
\newcommand{\Res}{\operatorname{Res}}
\newcommand{\Frac}{\operatorname{Frac}}
\newcommand{\End}{\operatorname{End}}
\newcommand{\Aut}{\operatorname{Aut}}
\newcommand{\supp}{\operatorname{supp}}
\newcommand{\trdeg}{\operatorname{trdeg}}
\newcommand{\id}{\operatorname{id}}
\newcommand{\calF}{\mathcal F}
\newcommand{\calR}{\mathcal R}
\newcommand{\calS}{\mathcal S}
\newcommand{\calW}{\mathcal W}
\newcommand{\calV}{\mathcal V}
\newcommand{\calB}{\mathcal B}
\newcommand{\calQ}{\mathcal Q}
\newcommand{\calL}{\mathcal L}
\newcommand{\qh}{\hbar}
\newcommand{\shuffle}{\mathbin{\amalg}}
\newcommand{\ev}{\operatorname{ev}}
\newcommand{\emptyword}{\varnothing}
\newcommand{\Li}{\operatorname{Li}}
\title{Rational-Exponential Iterated Integrals\\
and Quantum Polylogarithms\\[0.4em]
\large Absolute fixed-parameter completeness and a pentagon crossed product}
\author{Christopher D. Long}
\date{October 1, 2026\\\small Research draft, revised references}

\begin{document}
\maketitle
\begin{center}\small\texttt{galizur@gmail.com}\end{center}

\begin{abstract}
We study Goncharov's quantum polylogarithms at a fixed positive irrational
parameter. A Laplace transform of the cumulative momentum denominators gives
an ordered iterated-integral presentation, with the upper $i0$ prescription
justified on positively shifted contours. We compute the constant-linear
space of translated weight-zero Fourier letters as a two-variable localization
modulo Laurent polynomials. A residue argument makes this alphabet independent
modulo meromorphically exact differentials, yielding shuffle completeness at
all positive integral weights and depths in one common translation variable.
We also present the entire weight-zero differential field. After adjoining
$e^x$ and $e^{x/\hbar}$, it is a rational function field in the derivatives of
explicit normalized mixed towers. These derivatives are algebraically
independent: any dependence would produce an additive difference telescoper,
which a single-pole-orbit obstruction excludes. Combining the two presentations
classifies nonlinear coefficient relations and positive-weight relations over
$\C(x)$, without an unspecified meromorphic coefficient field. A separate
Ore-denominator argument proves the faithful cyclic crossed-product
presentation of Goncharov's pentagon operator algebra. The scalar and operator
statements remain separate. No parameter-changing completeness, scalar
quantum-cluster action, or arithmetic specialization theorem is asserted.
\end{abstract}

\medskip
\noindent\textit{Keywords.} Quantum polylogarithms; rational-exponential integrals;
iterated integrals; differential algebraic independence; difference equations;
shuffle algebra; quantum dilogarithm; Ore domain.

\clearpage
\tableofcontents
\clearpage

\section{Introduction and principal statements}\label{sec:intro}

Goncharov defines quantum polylogarithms by multiple Fourier integrals with
cumulative momentum denominators and proves modular difference equations,
continuation properties, and shuffle identities \cite{Goncharov2026}.
Their integrands are rational expressions in the momentum variables and their
exponentials. We use his term \emph{rational-exponential integrals} for this
analytic class. Our first purpose is to exhibit a direct ordered-integral
presentation of the positive-integral-weight sector. Our second is to determine
the constant-linear letter relations, the nonlinear algebraic structure of
the weight-zero differential field, and the polynomial relations among the
resulting iterated integrals.

Throughout the completeness theorems, $\qh$ and finitely many offsets are
fixed, and the variable is a common additive translation of the arguments.
The shuffle theorem is first proved over the actual single-valued meromorphic
weight-zero field. An independent difference-algebra argument then presents
that field as a rational function field in an explicit differential
transcendence basis. The resulting combined presentation is \emph{absolute}
in the sense that its coefficient field over $\C(x)$ is fully presented.
It is a functional relation theorem, not a numerical Kontsevich--Zagier
theorem, and uses no arithmetic specialization or $E$-function value lifting.

The analytic reduction is related to Chen's iterated integrals
\cite{Chen1977} and to the modified reduced-bar formalism used by Brown
\cite[Section~3]{Brown2009}. The general shuffle machinery is not new.
The specific inputs here are the contour-level reduction, the quantum-letter
quotient, the residue obstruction for iterated integrals, and a normalized
tower presentation of the nonlinear coefficient field. The latter uses the
classical connection between additive difference equations, telescoping,
and algebraic independence \cite{Schneider2010,HardouinSinger2008}; the
needed obstruction is proved directly for these towers.
In a separate appendix we give an operator result for the pentagon structure
of \cite{Goncharov2007}. It neither uses nor implies an action on the scalar
iterated-integral algebra.

\subsection{Concurrent and related developments}

Other $q$-deformations of iterated integrals address related but different
questions. Seki proves Hirose's duality conjecture for a $q$-discretization
on the four-punctured projective line, with word-dependent parameter shifts,
using a terminating basic-hypergeometric connector \cite{Seki2026}.
Bl\"umlein--Gavrilik--Mykhailiv--Schneider construct $q$-extensions of
iterated-integral alphabets and nested sums motivated by quantum field
theory \cite{BGMS2026}. We cite these as distinct deformation theories;
we neither identify their functions with Goncharov's Fourier integrals
nor use their word identities to obtain the completeness theorem here.

The contemporaneous double-sine constructions of cyclic Haagerup--Izumi
categories by Huang \cite{Huang2026} and Gannon--Schopieray--Yadav
\cite{GSY2026} belong to the finite categorical side of this subject.
They provide important context for the arithmetic direction, not inputs
to the fixed-parameter residue, coefficient-field, or shuffle proofs below.
The relation to RW and the Stark/SIC literature is specified in the
arithmetic-motivation subsection.

\subsection{The functions and the common translation variable}

Use the normalization
\[
 \sh z=e^z-e^{-z}=2\sinh z,
 \qquad
 H_{a,b}(p)=\frac1{\sh(\pi p)^a\sh(\pi\qh p)^b},
\]
where $a,b\in\N$ and $a+b>0$. The weight-zero Fourier letter is
\begin{equation}\label{eq:L-intro}
 L_{a,b}(z)=-\ii\int_{\R+\ii0}e^{-\ii pz}H_{a,b}(p)\dd p.
\end{equation}
Dependence on $\qh$ is suppressed. For $\mathbf n=(n_1,\ldots,n_m)$ with
$n_j\geq1$, put $|\mathbf n|=\sum n_j$ and
\begin{equation}\label{eq:F-intro}
 \calF_{\mathbf a,\mathbf b,\mathbf n}^{\qh}(\boldsymbol\omega)
 =\ii^{|\mathbf n|-m}
 \int_{(\R+\ii0)^m}
 \prod_{j=1}^m
 \frac{e^{-\ii p_j\omega_j}H_{a_j,b_j}(p_j)}
 {(p_1+\cdots+p_j)^{n_j}}
 \dd p_1\cdots\dd p_m .
\end{equation}
This is \cite[Definition~1.2]{Goncharov2026}, restricted to positive integral
weights and nontrivial kernels. The precise contour convention is given in
\cref{sec:contours}.

Fix a nonempty finite set $\Delta\subset\C$. Write
\begin{equation}\label{eq:translation}
 \omega_j=x+\delta_j-\ii\pi(r_j+s_j\qh),
 \qquad \delta_j\in\Delta,\quad r_j,s_j\in\Z.
\end{equation}
Set
\[
 \Lambda_{\qh}=\ii\pi(\Z+\qh\Z),\qquad
 \Delta/\Lambda_{\qh}
 =\{\text{cosets represented by elements of }\Delta\}.
\]
Let $\calW_{\qh,\Delta}$ be the complex vector space of actual functions spanned by
\begin{equation}\label{eq:alphabet}
 x\longmapsto L_{a,b}\bigl(x+\delta-\ii\pi(r+s\qh)\bigr).
\end{equation}
Let $K_{\qh,\Delta}$ be the differential subfield of $\Mer(\C)$ generated over
$\C(x)$ by these functions. Equivalently it is generated as a field by the
functions in \eqref{eq:alphabet} and all their derivatives. Its constants are
$\C$.

Define
\begin{equation}\label{eq:spectral-intro}
 \begin{split}
 \calR&=\C[X^{\pm1},Y^{\pm1}],\qquad
 d_X=X^{-1}-X,\quad d_Y=Y^{-1}-Y,\\
 \calS&=\calR[d_X^{-1},d_Y^{-1}],\qquad
 \calV_{\qh,\Delta}=\bigoplus_{\delta\in\Delta/\Lambda_{\qh}}\calS/\calR.
 \end{split}
\end{equation}
A representative $\delta$ is chosen in each coset. Different choices identify
the summands by multiplication by Laurent monomials.

\begin{theorem}[Letter presentation and residue exactness]\label{thm:mainletters}
Suppose $\qh>0$ is irrational. There is an isomorphism of complex vector spaces
\begin{equation}\label{eq:mainlettermap}
 \calV_{\qh,\Delta}\xrightarrow{\sim}\calW_{\qh,\Delta},\qquad
 \left[X^rY^s d_X^{-a}d_Y^{-b}\right]_{\delta}
 \longmapsto L_{a,b}\bigl(x+\delta-\ii\pi(r+s\qh)\bigr).
\end{equation}
The only constant-linear relations in one coset are the two centered
modular lowering relations, with the terminal symbol $e_{0,0}$ set equal to
zero. Moreover,
\begin{equation}\label{eq:mainexactness}
 \Omega\in\calW_{\qh,\Delta}\dd x,\quad
 \Omega=dG,\quad G\in\Mer(\C)
 \quad\Longrightarrow\quad \Omega=0.
\end{equation}
In particular $\calW_{\qh,\Delta}\dd x$ embeds into
$K_{\qh,\Delta}\dd x/dK_{\qh,\Delta}$.
\end{theorem}

For $f_1,\ldots,f_m\in\calW_{\qh,\Delta}$ define the ascending Chen integral
\begin{equation}\label{eq:Chen-intro}
 I[f_1|\cdots|f_m](x)
 =\int_{-\infty<t_1<\cdots<t_m<x}
 \prod_{j=1}^m f_j(t_j)\dd t_1\cdots\dd t_m.
\end{equation}
Initially all variables lie on the horizontal ray ending at $x$, in a
sufficiently far left half-plane. These integrals converge and give holomorphic
functions there. The empty word has value $1$.

\begin{theorem}[Shuffle completeness over the weight-zero field]\label{thm:mainshuffle}
Suppose $\qh>0$ is irrational and $\Delta\subset\C$ is finite and nonempty. Let
$\calQ_{\qh,\Delta}^{>0}$ be the $K_{\qh,\Delta}$-algebra generated by all
functions \eqref{eq:F-intro} on the common-translation lines
\eqref{eq:translation}, with arbitrary depth, $a_j,b_j\geq0$,
$a_j+b_j>0$, and $n_j\geq1$. Use the branch continued from the defining
strips into $\Re\omega_j<0$.
Then iterated integration induces an isomorphism
\begin{equation}\label{eq:mainshuffle}
 K_{\qh,\Delta}\otimes_{\C}\Sh(\calV_{\qh,\Delta})
 \xrightarrow{\sim}\calQ_{\qh,\Delta}^{>0}.
\end{equation}
The isomorphism uses \eqref{eq:mainlettermap} on letters. In particular, after
choosing an ordered basis of $\calV_{\qh,\Delta}$, the iterated integrals of
Lyndon words are algebraically independent over $K_{\qh,\Delta}$ and generate
the algebra. No generic-position hypothesis on the offsets is required:
coincident lattice cosets are combined as in \eqref{eq:spectral-intro}.
\end{theorem}

Here $\Sh(V)=\bigoplus_{n\geq0}V^{\otimes n}$ denotes the shuffle algebra,
not the tensor algebra with concatenation product. The theorem presents the
\emph{iterated-integral extension} of the coefficient field. It does not
identify $\calV_{\qh,\Delta}$ with the entire de~Rham quotient of that field.

\subsection{The coefficient field and the absolute presentation}

Choose a representative $\delta$ for each coset in $\Delta/\Lambda_\qh$, and
let
\begin{equation}\label{eq:representatives}
 \mathscr R_\Delta
 =\{\delta+\pi\ii(e+f\qh):\delta\in\Delta/\Lambda_\qh,\ e,f\in\{0,1\}\}.
\end{equation}
Here and below the notation $\delta\in\Delta/\Lambda_\qh$ refers to those
chosen representatives. Thus $\mathscr R_\Delta$ represents
$(\Delta+\Lambda_\qh)/(2\Lambda_\qh)$, with four representatives per original
coset. Put
\begin{equation}\label{eq:tower-intro}
 J_{b,\rho}(x)=L_{1,b}(x+\rho+\pi\ii\qh(b-1)),\qquad
 b\geq1,\quad\rho\in\mathscr R_\Delta,
 \qquad D=\frac{d}{dx}.
\end{equation}

\begin{theorem}[Explicit nonlinear coefficient field]\label{thm:maincoeff}
Fix $\qh>0$ irrational and finite nonempty $\Delta$. The functions
$U=e^x$ and $V=e^{x/\qh}$ are algebraically independent over $\C(x)$, and
\begin{equation}\label{eq:maincoeff}
 K_{\qh,\Delta}
 =\C(x,U,V)\bigl(D^jJ_{b,\rho}:b\geq1,\ j\geq0,
                                     \ \rho\in\mathscr R_\Delta\bigr).
\end{equation}
All the displayed derivatives are algebraically independent over
$\C(x,U,V)$. In particular the $J_{b,\rho}$ form a differential
transcendence basis and $K_{\qh,\Delta}$ is a purely differentially
transcendental extension of $\C(x,U,V)$.
\end{theorem}

Let $\mathscr A_\Delta$ be an ordered complex basis of $\calV_{\qh,\Delta}$,
and write $\mathscr L_\Delta$ for its Lyndon words. Introduce independent
indeterminates $u,v$, $y_{b,\rho,j}$ indexed as in \eqref{eq:maincoeff}, and
$t_\ell$ for $\ell\in\mathscr L_\Delta$. Every polynomial or rational
expression uses only finitely many of them.

\begin{corollary}[Absolute fixed-parameter functional completeness]
\label{cor:mainabsolute}
Evaluation gives a $\C(x)$-algebra isomorphism
\begin{equation}\label{eq:mainabsolute}
 \C(x)\bigl(u,v,y_{b,\rho,j}\bigr)
       [t_\ell:\ell\in\mathscr L_\Delta]
 \xrightarrow{\sim}\calQ_{\qh,\Delta}^{>0},
\end{equation}
where $u\mapsto e^x$, $v\mapsto e^{x/\qh}$,
$y_{b,\rho,j}\mapsto D^jJ_{b,\rho}$, and $t_\ell\mapsto I[\ell]$.
Consequently these functions are jointly algebraically independent over
$\C(x)$, and every element on the right reduces to a polynomial in the
Lyndon integrals with rational coefficients in the displayed independent
weight-zero variables.
\end{corollary}

The target in \eqref{eq:mainabsolute} includes the entire differential field
$K_{\qh,\Delta}$ by its definition in \cref{thm:mainshuffle}; it is not merely
the positive-weight algebra over $\C(x)$. Thus \cref{cor:mainabsolute}
presents the nonlinear coefficient field as well as its shuffle extension.
\Cref{thm:maincoeff} is proved in \cref{sec:coefficients}; the combined
statement is proved in \cref{sec:absolute-completeness}.

\subsection{What the proof establishes}

The scalar proof has four steps. First, the two tails of
\eqref{eq:L-intro} give a residue array with generating function
$-d_X^{-a}d_Y^{-b}$. Irrationality makes its lattice labels injective. Second,
the cumulative denominators in \eqref{eq:F-intro} are Laplace transformed into
powers of gaps in an ordered simplex. Polynomial gap factors reduce to
translated weight-zero letters by a Fourier integration-by-parts identity.
Third, a minimal-word-length argument gives independence of Chen words over
$K_{\qh,\Delta}$. Fourth, the two modular translations reduce the coefficient
field to normalized mixed towers. Additive algebraic dependence would force
a rational difference telescoper; a pole at a single position in an infinite
multiplicative orbit excludes every such telescoper. This proves the
coefficient-field presentation and removes the unspecified base field from
the combined completeness statement.

\Cref{sec:limits} records the limits of these statements. In particular,
continuation of a proved identity along a common path is not an independence
theorem after numerical specialization. Radchenko--Wheeler's algebraicity
results for arithmetic special values \cite{RadchenkoWheeler2026} are related
motivation, not an input to the proofs below. Their finite-scheme questions
are not considered here.

\section{Contours and meromorphic Fourier letters}\label{sec:contours}

Fix $\qh>0$ and write
\[
 \eta_{\qh}=\min(1,\qh^{-1}),\qquad
 c_{a,b}=\pi(a+b\qh).
\]
For $0<\eta<\eta_{\qh}$ let $C_\eta=\R+\ii\eta$, oriented from left to right.
No zero of either hyperbolic-sine factor lies in
$0<\Im p<\eta_{\qh}$.

\begin{lemma}[Fixed-height definition of the upper prescription]\label{lem:height}
For $|\Im z|<c_{a,b}$, the integral
\begin{equation}\label{eq:Lheight}
 -\ii\int_{C_\eta}e^{-\ii pz}H_{a,b}(p)\dd p
\end{equation}
is absolutely convergent, is independent of $\eta\in(0,\eta_{\qh})$, and is
holomorphic in $z$. It equals the contour with a small upper detour around
$p=0$. Thus it defines \eqref{eq:L-intro} without taking a nonuniform
integrand limit at the origin.
\end{lemma}
\begin{proof}
On each fixed horizontal line avoiding the zeros of the denominator,
\[
 |H_{a,b}(t+\ii\eta)|\leq C_\eta e^{-c_{a,b}|t|}.
\]
If $z=u+\ii v$, then
$|e^{-\ii(t+\ii\eta)z}|=e^{tv+\eta u}$.
On compact substrips the product is bounded by an integrable exponential,
locally uniformly in $z$. The same remains true after any number of
$z$-derivatives, since only powers of $p$ are introduced.
Cauchy's theorem on rectangles between two positive heights applies; the
vertical sides tend to zero by this estimate. Deformation to a path that is
real outside a small disk and passes above the origin crosses no pole.
\end{proof}

\subsection{A two-tail expansion}

For $a\in\N$ set
\begin{equation}\label{eq:Ca}
 C_a(m)=
 \begin{cases}
 \binom{a+m-1}{m},&a>0,\\
 1,&a=0,\ m=0,\\
 0,&a=0,\ m>0,
 \end{cases}
 \qquad m\in\N.
\end{equation}
For $p>0$ one has
\begin{equation}\label{eq:positiveexp}
 H_{a,b}(p)=\sum_{m,n\geq0}C_a(m)C_b(n)e^{-A_{m,n}p},\qquad
 A_{m,n}=\pi\bigl(a+2m+\qh(b+2n)\bigr).
\end{equation}
Terms with zero coefficient are omitted; every remaining $A_{m,n}$ is positive.

\begin{proposition}[Normally convergent meromorphic expansion]\label{prop:tails}
Choose $0<\rho<\eta_{\qh}$ and let $\Gamma_\rho$ be the upper semicircle from
$-\rho$ to $\rho$. Define the entire function
\[
 E_\rho(z)=-\ii\int_{\Gamma_\rho}e^{-\ii pz}H_{a,b}(p)\dd p.
\]
Then $L_{a,b}$ extends to a single-valued meromorphic function on $\C$, and
\begin{equation}\label{eq:tails}
 \begin{split}
 L_{a,b}(z)=E_\rho(z)-\ii\sum_{m,n\geq0}C_a(m)C_b(n)
 \left\{
 \frac{e^{-\rho(A_{m,n}+\ii z)}}{A_{m,n}+\ii z}
 +(-1)^{a+b}\frac{e^{-\rho(A_{m,n}-\ii z)}}{A_{m,n}-\ii z}
 \right\}.
 \end{split}
\end{equation}
Both series converge normally on compact sets away from their displayed
poles. If $\qh$ is irrational, all nonzero terms within either tail have
distinct poles, and
\begin{equation}\label{eq:residues}
 \Res_{z=\ii A_{m,n}}L_{a,b}(z)\dd z=-C_a(m)C_b(n),\qquad
 \Res_{z=-\ii A_{m,n}}L_{a,b}(z)\dd z=(-1)^{a+b}C_a(m)C_b(n).
\end{equation}
In particular every pole is simple.
\end{proposition}
\begin{proof}
Use \cref{lem:height} to split the contour into its two real tails and
$\Gamma_\rho$. On $[\rho,\infty)$ the expansion \eqref{eq:positiveexp} is
absolutely convergent. In the defining strip it may be integrated term by
term, giving the first series in \eqref{eq:tails}. On the negative tail put
$p=-q$ and use $H_{a,b}(-q)=(-1)^{a+b}H_{a,b}(q)$.

The coefficients $C_a(m)C_b(n)$ have at most polynomial growth, whereas
$e^{-\rho A_{m,n}}$ decays exponentially in every active index. On a fixed
compact set of $z$ values, all sufficiently large denominators have size
comparable to $A_{m,n}$. The finitely many other denominators are bounded
away from zero when the displayed poles are excluded. This proves normal
convergence. The compact contour contribution is entire by differentiation
under its integral.

The possible pole sets are locally finite, since $\qh>0$ and a bound on
$A_{m,n}$ bounds all active indices. Equality of two positive pole labels
implies $(m-m')+\qh(n-n')=0$, hence equality of the labels for irrational
$\qh$. The residues follow directly from the linear denominators in
\eqref{eq:tails}; their exponential numerators are $1$ at the poles.
\end{proof}

\begin{remark}
Although $\Z+\qh\Z$ is dense in $\R$ when $\qh$ is real irrational, the
\emph{positive cone} appearing in one letter's pole set is locally finite.
Only finite unions of translated cones are used in any argument. The
geometric expansions in \eqref{eq:positiveexp} are never applied across
$p=0$.
\end{remark}

\subsection{Lowering, multiplication, and reflection}

\begin{proposition}[Exact letter identities]\label{prop:identities}
For $a,b\geq0$, $a+b>0$, the meromorphic letters satisfy
\begin{align}
 L_{a,b}(z+\ii\pi)-L_{a,b}(z-\ii\pi)&=L_{a-1,b}(z)
 &&(a\geq1),\label{eq:lowerX}\\
 L_{a,b}(z+\ii\pi\qh)-L_{a,b}(z-\ii\pi\qh)&=L_{a,b-1}(z)
 &&(b\geq1),\label{eq:lowerY}
\end{align}
where the terminal symbol on the right is defined to be $L_{0,0}=0$.
Furthermore,
\begin{equation}\label{eq:mult}
 \begin{split}
 zL_{a,b}(z)
 ={}&\ii\pi a\bigl(L_{a+1,b}(z+\ii\pi)+L_{a+1,b}(z-\ii\pi)\bigr)\\
 &+\ii\pi\qh b\bigl(L_{a,b+1}(z+\ii\pi\qh)
                    +L_{a,b+1}(z-\ii\pi\qh)\bigr).
 \end{split}
\end{equation}
Finally,
\begin{equation}\label{eq:reflection}
 L_{a,b}(z)-(-1)^{a+b}L_{a,b}(-z)=P_{a,b}(z),\qquad
 P_{a,b}(z)=-2\pi\Res_{p=0}\bigl(e^{-\ii pz}H_{a,b}(p)\bigr),
\end{equation}
where $P_{a,b}$ is a polynomial of degree $a+b-1$.
\end{proposition}
\begin{proof}
If the lowered index is nonterminal, subtract the integrands on a common
convergence strip. The multipliers are $e^{\pi p}-e^{-\pi p}$ and
$e^{\pi\qh p}-e^{-\pi\qh p}$, respectively. This proves
\eqref{eq:lowerX}--\eqref{eq:lowerY} on a nonempty strip, and then everywhere
by meromorphic continuation. These are also the depth-zero cases of
\cite[Section~3.0.1]{Goncharov2026}.

The terminal cases require a separate check; they are not integrals of a
constant Fourier multiplier. Closing a rectangular contour upwards for
$\Re z<0$ gives
\begin{equation}\label{eq:axesbase}
 L_{1,0}(z)=\sum_{k\geq1}(-1)^ke^{kz}
          =-\frac{e^z}{1+e^z},\qquad
 L_{0,1}(z)=\qh^{-1}L_{1,0}(z/\qh).
\end{equation}
For the first identity the residues at $p=\ii k$, $k\geq1$, are
$(-1)^ke^{kz}/(2\pi)$; heights $k+1/2$ make the top edge tend to zero.
The second follows by rescaling $p$. The required terminal centered
differences vanish by the respective $2\pi\ii$ and $2\pi\ii\qh$
periodicities. This justifies $L_{0,0}=0$ solely as a notation for the terminal
\emph{identity}.

For \eqref{eq:mult}, integrate
$\partial_p(e^{-\ii pz}H_{a,b}(p))$ along $C_\eta$.
The endpoints vanish exponentially in the defining strip. Since
\[
 H'_{a,b}(p)=-\pi a(e^{\pi p}+e^{-\pi p})H_{a+1,b}(p)
             -\pi\qh b(e^{\pi\qh p}+e^{-\pi\qh p})H_{a,b+1}(p),
\]
and $\int e^{-\ii pz}H_{a,b}(p)\dd p=\ii L_{a,b}(z)$, the result is
\eqref{eq:mult}. Every shifted integral converges in the original strip.

For reflection, substitute $p\mapsto-p$ in the reflected integral. The
upper contour becomes the lower contour, and the parity of $H_{a,b}$ is
$(-1)^{a+b}$. The upper contour minus the lower one is clockwise around the
origin, so multiplication by $-\ii$ gives the factor $-2\pi$ in
\eqref{eq:reflection}. The pole order at the origin is $a+b$, and the residue
of its product with $e^{-\ii pz}$ is a polynomial of the stated degree.
\end{proof}

\begin{remark}
The reflection correction is a polynomial, hence its product with $\dd z$
is exact over $\C(z)$. It must not be discarded as a zero \emph{function}.
Our residue normalization gives $P_{1,0}=-1$, consistently with
\eqref{eq:axesbase}. Compare the quantum Bernoulli reflection formulas in
\cite[Section~5]{Goncharov2026}.
\end{remark}

\section{Residue tails and the complete letter quotient}\label{sec:residues}

Assume from now until \cref{sec:laplace} that $\qh>0$ is irrational. Put
$T_{r,s}f(z)=f(z-\ii\pi(r+s\qh))$. Use the rings
$\calR,\calS$ and elements $d_X,d_Y$ of \eqref{eq:spectral-intro}.

\subsection{The algebraic presentation}

\begin{lemma}[Localization presentation]\label{lem:presentation}
Let $M$ be the $\calR$-module generated by $e_{a,b}$, $a,b\geq0$, with
\begin{equation}\label{eq:module-rel}
 e_{0,0}=0,\qquad d_Xe_{a,b}=e_{a-1,b}\ (a\geq1),\qquad
 d_Ye_{a,b}=e_{a,b-1}\ (b\geq1).
\end{equation}
Then
\begin{equation}\label{eq:module-iso}
 M\xrightarrow{\sim}\calS/\calR,\qquad
 e_{a,b}\longmapsto[d_X^{-a}d_Y^{-b}].
\end{equation}
\end{lemma}
\begin{proof}
The displayed map respects every relation and is surjective. For injectivity,
write an arbitrary finite element as $\sum P_{a,b}e_{a,b}$ and choose $A,B$
bounding its indices. The relations give
\[
 \sum P_{a,b}e_{a,b}=Qe_{A,B},\qquad
 Q=\sum P_{a,b}d_X^{A-a}d_Y^{B-b}.
\]
If its image in $\calS/\calR$ vanishes, then
$Q=d_X^Ad_Y^B P$ for some $P\in\calR$. Thus
$Qe_{A,B}=Pe_{0,0}=0$.
\end{proof}

By \cref{prop:identities}, $X,Y$ may act on actual letters by $T_{1,0}$ and
$T_{0,1}$; \eqref{eq:module-rel} then gives a well-defined map from $M$
to their constant-linear span. We prove that this analytic map is injective.

\subsection{Tail injectivity}

There is an injective expansion map
\begin{equation}\label{eq:expansionmap}
 \calS\hookrightarrow\C[[X,Y]][X^{-1},Y^{-1}],\qquad
 d_X^{-a}d_Y^{-b}\longmapsto
 \sum_{m,n\geq0}C_a(m)C_b(n)X^{a+2m}Y^{b+2n}.
\end{equation}
Indeed $d_X=(1-X^2)/X$, and similarly for $Y$; denominators $1-X^2$ and
$1-Y^2$ are units in the formal power-series ring. The expansion has finite
support if and only if the original element lies in $\calR$. Consequently
\begin{equation}\label{eq:tail-inj}
 \calS/\calR\hookrightarrow
 \C[[X,Y]][X^{-1},Y^{-1}]/\calR
\end{equation}
is injective.

\begin{theorem}[Residue-tail detection]\label{thm:tail}
For a finite expression
\begin{equation}\label{eq:finiteOmega}
 \Omega(z)=\sum c_{a,b,r,s}\,T_{r,s}L_{a,b}(z)\dd z
\end{equation}
set
\[
 F_\Omega(X,Y)=\sum c_{a,b,r,s}X^rY^sd_X^{-a}d_Y^{-b}\in\calS.
\]
If all finite residues of $\Omega$ vanish, then $F_\Omega\in\calR$.
In particular, a nonzero class $[F_\Omega]\in\calS/\calR$ forces a nonzero
finite residue of $\Omega$.
\end{theorem}
\begin{proof}
For a translated letter the positive poles from \cref{prop:tails} have labels
\begin{equation}\label{eq:positive-label}
 (k,\ell)=(r+a+2m,s+b+2n),\qquad
 z=\ii\pi(k+\qh\ell),
\end{equation}
and residues $-C_a(m)C_b(n)$. Their generating series is therefore
$-X^rY^sd_X^{-a}d_Y^{-b}$.

The negative poles have labels
$(r-a-2m,s-b-2n)$. A finite collection of positive supports lies in a finite
union of coordinatewise lower-bounded quadrants, and a finite collection of
negative supports lies in a finite union of coordinatewise upper-bounded
quadrants. Their intersection is finite. Since
$(k,\ell)\mapsto k+\qh\ell$ is injective on $\Z^2$, positive and negative
poles can therefore cancel at only finitely many labels.

If all actual residues vanish, the expansion of $F_\Omega$ has finite
support. The injectivity statement \eqref{eq:tail-inj} implies
$F_\Omega\in\calR$.
\end{proof}

\begin{corollary}[One-coset exactness]\label{cor:onecoset}
The map
\[
 \calS/\calR\longrightarrow
 \operatorname{span}_{\C}\{T_{r,s}L_{a,b}\},\qquad
 [X^rY^sd_X^{-a}d_Y^{-b}]\longmapsto T_{r,s}L_{a,b}
\]
is an isomorphism. A form in this span is the differential of a globally
single-valued meromorphic function if and only if it is the zero form.
\end{corollary}
\begin{proof}
Surjectivity and well-definedness follow from \cref{lem:presentation} and
\cref{prop:identities}. A zero analytic expression has zero residues, so
\cref{thm:tail} makes its formal class zero; this proves injectivity.
If $G$ is meromorphic, its Laurent derivative has no residue at any finite
point. Apply \cref{thm:tail} and then the same injectivity argument.
The converse is immediate.
\end{proof}

\begin{remark}\label{rem:notcontact}
The quotient by $\calR$ in \eqref{eq:module-iso} removes \emph{finite-support
residue tails}. It is not an identification of Laurent multipliers with
contact distributions. The passage from a zero tail class to an actual zero
function uses the exact lowering identities, including the separately verified
terminal identities. Residues by themselves would not exclude a nonzero
entire remainder.
\end{remark}

\subsection{Offsets and an explicit basis}

The possible poles of letters in the $\delta$-coset lie in
$-\delta+\Lambda_{\qh}$. For inequivalent offsets these sets are disjoint.
For equivalent offsets, their letters are identified by the corresponding
Laurent monomial. Applying \cref{cor:onecoset} in each disjoint coset proves
\cref{thm:mainletters}.

There is also a simple vector-space basis for the quotient, so the statement
does not depend on an unspecified cohomology calculation.

\begin{proposition}[Partial-fraction basis]\label{prop:basis}
A complex basis of $\calS/\calR$ is given by the classes of
\begin{align}
 &(X-\epsilon)^{-a}Y^n,
 &&\epsilon\in\{1,-1\},\ a\geq1,\ n\in\Z,\label{eq:basis1}\\
 &X^m(Y-\epsilon)^{-b},
 &&\epsilon\in\{1,-1\},\ b\geq1,\ m\in\Z,\label{eq:basis2}\\
 &(X-\epsilon)^{-a}(Y-\epsilon')^{-b},
 &&\epsilon,\epsilon'\in\{1,-1\},\ a,b\geq1.\label{eq:basis3}
\end{align}
Every basis class is represented by a finite constant-linear combination of
translated Fourier letters.
\end{proposition}
\begin{proof}
Partial fractions in one variable give the direct sum
\[
 \C[X^{\pm1},(1-X^2)^{-1}]
 =\C[X^{\pm1}]
 \oplus\bigoplus_{\epsilon=\pm1,\ a\geq1}\C(X-\epsilon)^{-a}.
\]
Tensor this decomposition with its $Y$-analogue and quotient out the
Laurent--Laurent summand. This gives precisely
\eqref{eq:basis1}--\eqref{eq:basis3} and their independence.
Finally
\[
 \frac1{X-\epsilon}=-\frac{X+\epsilon}{Xd_X}
\]
expresses every such basis element in the presentation of
\cref{lem:presentation}; the assertion about letters follows.
\end{proof}

\section{Left-half-plane normalization and polynomial closure}\label{sec:decay}

The residue results are global, whereas the iterated integrals need a
compatible normalization at an infinite endpoint. We establish it before
using any bar construction.

\begin{lemma}[Decay of all fixed translates]\label{lem:decay}
Let $\qh>0$, and fix $a,b$ with $a+b>0$. For every compact interval
$J\subset\R$ and every $0<\eta<\eta_{\qh}$,
\begin{equation}\label{eq:decay}
 L_{a,b}(u+\ii v)=O(e^{\eta u})\quad(u\longrightarrow-\infty),
 \qquad v\in J.
\end{equation}
The estimate is locally uniform for every fixed finite family of translates.
Each function in $\calW_{\qh,\Delta}$ is holomorphic in a sufficiently far
left half-plane and satisfies such an estimate there.
\end{lemma}
\begin{proof}
On a compact substrip of $|\Im z|<c_{a,b}$, \cref{lem:height} on the line
$C_\eta$ gives \eqref{eq:decay} directly.
For the one-frequency kernels, \eqref{eq:mult} and
\eqref{eq:lowerX} give
\begin{equation}\label{eq:axisrecursion}
 L_{a+1,0}(z)=\frac{z+\ii\pi(a-1)}{2\pi\ii a}\,
 L_{a,0}(z-\ii\pi)\qquad(a\geq1).
\end{equation}
Together with \eqref{eq:axesbase}, this expresses $L_{a,0}$ as a polynomial
of degree $a-1$ times a shifted logistic function. It therefore satisfies
\eqref{eq:decay} on any bounded imaginary interval for $\eta<1$.
Rescaling gives $L_{0,b}(z)=\qh^{-1}L_{b,0}(z/\qh)$ and the assertion
for $a=0$.

For $a,b>0$, rewrite \eqref{eq:lowerX} as
\[
 L_{a,b}(z+2\pi\ii)=L_{a,b}(z)+L_{a-1,b}(z+\pi\ii).
\]
Since $c_{a,b}>\pi$, any fixed bounded imaginary interval can be covered by
finitely many translates of compact substrips of the defining strip by
$2\pi\ii\Z$. Iterating the displayed identity reduces its additional terms
to the case with smaller $a$. Induction, starting at $a=0$, proves the
estimate. All poles lie on the imaginary axis by \cref{prop:tails}; after
finitely many offsets a left half-plane avoids them all.
\end{proof}

\begin{proposition}[Polynomial multiples remain in the letter space]\label{prop:polyclosure}
For every finite offset set $\Delta$ one has
\begin{equation}\label{eq:polyclosure}
 \C[x]\calW_{\qh,\Delta}\subseteq\calW_{\qh,\Delta}.
\end{equation}
More explicitly, any polynomial multiple of a letter in
\eqref{eq:alphabet} is a finite constant-linear combination of letters in
its same offset coset.
\end{proposition}
\begin{proof}
For $u=x+\delta-\ii\pi(r+s\qh)$, write $x=u-\delta+\ii\pi(r+s\qh)$.
Formula \eqref{eq:mult} expresses $uL_{a,b}(u)$ in translated letters.
Iteration proves the result for powers of $x$, and then for all polynomials.
\end{proof}

\begin{remark}[Reflected letters]
Reflection need not preserve the decay normalization at $-\infty$.
By \eqref{eq:reflection}, adjoining reflected letters, and if necessary
replacing $\Delta$ by $\Delta\cup(-\Delta)$, enlarges the alphabet only by
polynomials. Their differentials are exact. In fact
$\calW_{\qh,\Delta}\cap\C[x]=0$ by \cref{lem:decay}, and
\cref{thm:mainletters} implies that an element of
$(\calW_{\qh,\Delta}+\C[x])\dd x$ is meromorphically exact precisely when
its $\calW$-component is zero. Arbitrary reflected words at an infinite
endpoint require boundary subtraction or a finite basepoint; no such
regularization is implicit in our main theorem.
\end{remark}

\section{The nonlinear weight-zero field}\label{sec:coefficients}

We now present the coefficient field itself. The residue theorem of
\cref{sec:residues} detects constant-linear exactness, but does not by itself
exclude polynomial relations between meromorphic functions. For that purpose
we use an additive difference equation and a rational pole-orbit obstruction.
The general relation between telescoping and algebraic independence is
classical; see \cite{Schneider2010,HardouinSinger2008}. The particular lemmas
needed below are proved directly. No differential closure of the difference
constants, or assertion that an analytic extension has no new difference
constants, is required.

Throughout this section $\qh>0$ is irrational, $\Delta$ is finite and nonempty,
and $D=d/dx$. Use the representative set $\mathscr R_\Delta$ and the towers
$J_{b,\rho}$ of \eqref{eq:tower-intro}. Write
\begin{equation}\label{eq:elementarybase}
 U=e^x,\qquad V=e^{x/\qh},\qquad F_\qh=\C(x,U,V),\qquad
 \mathfrak q=e^{2\pi\ii/\qh}.
\end{equation}
Here $U,V$ denote functions; the proof below identifies them with independent
rational-function variables over $\C(x)$. The parameter $\mathfrak q$ in this
section is distinct from the quantum-torus parameter in \cref{app:operator}.
Define commuting automorphisms on $\Mer(\C)$ by
\[
 (\sigma f)(x)=f(x+2\pi\ii),\qquad
 (\tau f)(x)=f(x+2\pi\ii\qh).
\]
They commute with $D$ and preserve $F_\qh$.

\subsection{The elementary field and its difference constants}

\begin{lemma}[Elementary independence and fixed field]\label{lem:fixedfield}
The functions $U,V$ are algebraically independent over $\C(x)$. On their
rational function field,
\begin{equation}\label{eq:sigma-base}
 \sigma(x)=x+2\pi\ii,\quad \sigma(U)=U,\quad
 \sigma(V)=\mathfrak qV,\qquad F_\qh^\sigma=\C(U).
\end{equation}
Furthermore,
\begin{equation}\label{eq:D-base}
 D=\partial_x+U\partial_U+\qh^{-1}V\partial_V
 \quad\hbox{on }F_\qh.
\end{equation}
\end{lemma}
\begin{proof}
After clearing denominators, a polynomial relation between $U$ and $V$ would
be a finite identity
\[
 \sum_{m,n}p_{m,n}(x)e^{(m+n/\qh)x}=0,
 \qquad p_{m,n}\in\C[x].
\]
Distinct pairs $(m,n)$ give distinct real rates because $\qh$ is irrational.
Divide by the exponential with largest rate and let real $x\to+\infty$.
All terms but its polynomial coefficient tend to zero. A nonzero polynomial
cannot tend to zero in this way. Removing that term and repeating proves
independence.

Put $E=\C(U)$, $t=2\pi\ii$, and regard $F_\qh=E(x,V)$. The automorphism
fixes $E$ and sends $(x,V)$ to $(x+t,\mathfrak qV)$.
Suppose $0\ne f=P/Q\in F_\qh$ is fixed, where $P,Q\in E[x,V]$ are coprime.
The identity $\sigma(P)Q=P\sigma(Q)$ implies that $P$ divides $\sigma(P)$.
Total degree is preserved, so $\sigma(P)=cP$ with $c\in E^\times$;
then also $\sigma(Q)=cQ$.

Write $P=\sum_np_n(x)V^n$. For every nonzero coefficient polynomial,
\[
 \mathfrak q^n p_n(x+t)=c p_n(x).
\]
Leading coefficients in $x$ give $c=\mathfrak q^n$. As $\mathfrak q$ is not
a root of unity, only one $n$ can occur. The remaining $p_n$ is a periodic
polynomial, hence belongs to $E$. Thus $P=aV^n$. The same argument gives
$Q=bV^m$, and the common multiplier forces $n=m$. Hence $f=a/b\in E$.
This proves \eqref{eq:sigma-base}; \eqref{eq:D-base} is the chain rule.
\end{proof}

\begin{remark}
The difference-constant field in \eqref{eq:sigma-base} is $\C(U)$, not $\C$.
In particular the telescoping coefficients below must be allowed to be
arbitrary rational functions of $U$. By contrast, the constants for the
ordinary derivative $D$ in every field between $\C$ and $\Mer(\C)$ are
$\C$. The two notions of constants must not be interchanged.
\end{remark}

\subsection{Additive dependence implies a rational telescoper}

\begin{lemma}[Additive-dependence obstruction]\label{lem:additive-dependence}
Let $F\subseteq E$ be difference fields of characteristic zero, with the
same notation $\sigma$ for their compatible automorphisms. Suppose
$y_1,\ldots,y_m\in E$ satisfy
\[
 \sigma(y_i)=y_i+b_i,\qquad b_i\in F.
\]
If the $y_i$ are algebraically dependent over $F$, then there are
$c_i\in F^\sigma$, not all zero, and $g\in F$ such that
\begin{equation}\label{eq:additive-telescoper}
 \sum_i c_i b_i=\sigma(g)-g.
\end{equation}
No hypothesis on $E^\sigma$ is made.
\end{lemma}
\begin{proof}
Let $\mathfrak a$ be the nonzero kernel of evaluation
$F[Z_1,\ldots,Z_m]\to E$, $Z_i\mapsto y_i$.
The semilinear affine operator
\[
 \mathcal T(P)=\sigma(P)(Z_1+b_1,\ldots,Z_m+b_m)
\]
preserves $\mathfrak a$, since evaluation of $\mathcal T(P)$ is
$\sigma(P(y_1,\ldots,y_m))$.
Choose a graded monomial order and a nonzero $P\in\mathfrak a$ whose leading
monomial is least possible, with leading coefficient normalized to $1$.
Translation preserves its leading monomial and coefficient. Therefore
$\mathcal T(P)-P$, if nonzero, would contradict minimality. Thus
$\mathcal T(P)=P$.

Let $n=\deg_ZP\geq1$. An appropriate formal partial derivative of total order
$n-1$ is a nonconstant affine-linear polynomial
\[
 Q(Z)=c_0+\sum_i c_iZ_i.
\]
Formal partial derivatives in the $Z_i$ commute with $\mathcal T$, because
$b_i\in F$. Hence $\mathcal T(Q)=Q$. Comparing coefficients gives
$\sigma(c_i)=c_i$ and
$\sigma(c_0)+\sum_i c_i b_i=c_0$.
At least one $c_i$ is nonzero, and $g=-c_0$ proves the claim.
The derivative $Q$ need not itself lie in $\mathfrak a$; its invariance is
all that is used.
\end{proof}

We use only the implication in this lemma. The existence of a telescoper
need not imply algebraic dependence if the extension introduces new
$\sigma$-constants. This is why no converse, or unchanged-constants assertion,
is implicit here. Applying the lemma to finitely many derivatives of the
$y_i$ is legitimate whenever $D\sigma=\sigma D$.

\subsection{A pole at a single orbit position cannot telescope}

In this subsection, poles mean valuations of the rational function field
$E(x)(V)$ along the linear polynomials $V-\alpha$, with $\alpha\in\C^\times$.
They are algebraic poles in the independent variable $V$, not assertions
about a topological ordering of poles of a meromorphic function of $x$.

\begin{lemma}[Finite-string obstruction]\label{lem:orbit-poles}
Let $E$ be a field containing $\C$, let $t\in\C^\times$, and let
$q\in\C^\times$ have infinite multiplicative order. On $E(x,V)$ let
$\sigma$ fix $E$ and send $(x,V)$ to $(x+t,qV)$.
If $r=\sigma(g)-g$ with $g\in E(x,V)$, then, in each orbit
$\{\alpha q^n:n\in\Z\}$, the set of finite $V$-poles of $r$ is either
empty or contains at least two positions.
\end{lemma}
\begin{proof}
For a nonzero $g$, the set
\[
 S=\{n\in\Z:g\text{ has a pole along }V-\alpha q^n\}
\]
is finite. The corresponding set for $\sigma(g)$ is $S-1$.
Indeed, substitution by $(x+t,qV)$ transports the local valuation at one
position to the preceding position and sends every nonzero leading
coefficient in $E(x)$ to a nonzero one.
If $S$ is empty, so is the pole set of $\sigma(g)-g$ on the orbit.
Otherwise let $m=\min S$, $n=\max S$. At $m-1$ only $\sigma(g)$ has a pole;
at $n$ only $g$ has a pole. Both poles survive in the difference, and
$m-1<n$.
An irreducible divisor not in this orbit cannot map into it under $\sigma$,
since the inverse image of $V-\alpha q^k$ is again a constant linear divisor
in the same orbit. Thus other, possibly $x$-dependent denominator factors
of $g$ do not affect this argument.
\end{proof}

\subsection{Normalized mixed towers and their increments}

For $b\geq1$ and $\rho\in\mathscr R_\Delta$, put
\begin{align}
 a_\rho&=\exp((\rho+\pi\ii)/\qh),&
 \alpha_\rho&=-a_\rho^{-1},&
 R_\rho(V)&=-\frac{a_\rho V}{1+a_\rho V}
           =-\frac{V}{V-\alpha_\rho},\label{eq:rho-logistic}\\
 A_{b,\rho}(x)&=
 \frac{\displaystyle\prod_{\ell=1}^{b-1}
 (x+\rho+\pi\ii+\pi\ii\qh(2\ell-1))}
 {\qh(2\pi\ii\qh)^{b-1}(b-1)!}.&&&\label{eq:tower-polynomial}
\end{align}
The empty product gives $A_{1,\rho}=1/\qh$.

\begin{proposition}[Tower difference equations]\label{prop:tower-increments}
The towers satisfy
\begin{equation}\label{eq:tower-sigma}
 \sigma J_{b,\rho}-J_{b,\rho}=A_{b,\rho}(x)R_\rho(V).
\end{equation}
For $b\geq1$ they also satisfy
\begin{equation}\label{eq:tower-tau}
 \tau^{-1}J_{b,\rho}=J_{b,\rho}-J_{b-1,\rho},
 \qquad J_{0,\rho}(x):=L_{1,0}(x+\rho-\pi\ii\qh)\in F_\qh.
\end{equation}
The poles $\alpha_\rho$ lie in pairwise distinct $\mathfrak q$-orbits, and
$\deg_xA_{b,\rho}=b-1$.
\end{proposition}
\begin{proof}
The first lowering identity gives
\[
 \sigma J_{b,\rho}-J_{b,\rho}
 =L_{0,b}(x+\rho+\pi\ii+\pi\ii\qh(b-1)).
\]
Iterating \eqref{eq:axisrecursion} gives the explicit formula
\begin{equation}\label{eq:axis-product}
 L_{b,0}(z)=
 \frac{\displaystyle\prod_{\ell=1}^{b-1}(z+\pi\ii(2\ell-b))}
 {(2\pi\ii)^{b-1}(b-1)!}
 L_{1,0}(z-\pi\ii(b-1)).
\end{equation}
Use $L_{0,b}(z)=\qh^{-1}L_{b,0}(z/\qh)$ and substitute the preceding
argument. The polynomial factor is \eqref{eq:tower-polynomial}, and the
remaining logistic factor is \eqref{eq:rho-logistic}. This proves
\eqref{eq:tower-sigma}. The second lowering identity, centered at
$x+\rho+\pi\ii\qh(b-2)$, gives \eqref{eq:tower-tau}, including $b=1$.

Finally,
\[
 \frac{\alpha_\rho}{\alpha_{\rho'}}\in\mathfrak q^{\Z}
 \quad\Longleftrightarrow\quad
 \rho-\rho'\in2\pi\ii(\Z+\qh\Z)=2\Lambda_\qh.
\]
This follows by using the kernel $2\pi\ii\Z$ of the complex exponential.
The representatives are distinct modulo $2\Lambda_\qh$. The degree assertion
is immediate from the displayed product.
\end{proof}

\subsection{Independence of every finite collection of tower derivatives}

\begin{proposition}[No rational telescoper for the tower increments]
\label{prop:no-telescope}
If finitely many $c_{b,\rho,j}(U)\in\C(U)$ and a function $g\in F_\qh$
satisfy
\begin{equation}\label{eq:jet-telescoper}
 \sum_{b,\rho,j}c_{b,\rho,j}(U)
 D^j\bigl(A_{b,\rho}(x)R_\rho(V)\bigr)=\sigma(g)-g,
\end{equation}
then every $c_{b,\rho,j}$ is zero.
\end{proposition}
\begin{proof}
Fix an active $\rho$ and let $j$ be the largest derivative order with a
nonzero coefficient for this $\rho$. At the valuation $V=\alpha_\rho$,
\begin{equation}\label{eq:leading-jet-pole}
 D^j(A_{b,\rho}R_\rho)
 =\frac{(-1)^{j+1}j!\,\alpha_\rho^{j+1}\qh^{-j}A_{b,\rho}(x)}
 {(V-\alpha_\rho)^{j+1}}
 +\text{terms of pole order at most }j.
\end{equation}
For $j=0$ this follows from
$R_\rho=-1-\alpha_\rho/(V-\alpha_\rho)$. Inductively, only the
$(V/\qh)\partial_V$ part of $D$ raises the pole order; its leading value at
the divisor is $\alpha_\rho/\qh$. Derivatives of the polynomial in $x$ give
lower pole orders.

The leading coefficient in the $\rho$-part of the left side of
\eqref{eq:jet-telescoper} is therefore a fixed nonzero scalar times
\[
 \sum_b c_{b,\rho,j}(U)A_{b,\rho}(x).
\]
It is nonzero: the $A_{b,\rho}$ have different degrees in the independent
variable $x$, so they are linearly independent over $\C(U)$.
Hence that part has a pole at $\alpha_\rho$, and nowhere else in its
$\mathfrak q$-orbit. Other $\rho'$ contribute no pole on the same orbit by
\cref{prop:tower-increments}. Thus the entire left side has exactly one pole
position on this orbit, contradicting \cref{lem:orbit-poles} with
$E=\C(U)$, $t=2\pi\ii$, and $q=\mathfrak q$.
\end{proof}

\begin{corollary}[Differential independence of the mixed towers]
\label{cor:tower-independent}
The functions $D^jJ_{b,\rho}$, for $b\geq1$, $j\geq0$, and
$\rho\in\mathscr R_\Delta$, are algebraically independent over $F_\qh$.
Equivalently the $J_{b,\rho}$ are differentially algebraically independent
over that differential field.
\end{corollary}
\begin{proof}
Apply \cref{lem:additive-dependence} to any finite collection of these
functions in the difference field $\Mer(\C)$. Their increments are
$D^j(A_{b,\rho}R_\rho)$ because $D$ and $\sigma$ commute. By
\cref{lem:fixedfield} the resulting coefficients would lie in $\C(U)$.
\Cref{prop:no-telescope} excludes the nontrivial telescoper, so the finite
collection is independent. Algebraic or differential-algebraic dependence
would involve only a finite collection.
\end{proof}

\subsection{Reduction of every weight-zero generator}

\begin{proposition}[Generation by the normalized towers]
\label{prop:tower-generate}
One has
\begin{equation}\label{eq:tower-generation}
 K_{\qh,\Delta}
 =F_\qh\bigl(D^jJ_{b,\rho}:b\geq1,\ j\geq0,
                         \ \rho\in\mathscr R_\Delta\bigr).
\end{equation}
Every original translated Fourier letter and every derivative has a finite
rational expression in the displayed generators. Such a reduction uses only
\eqref{eq:axesbase}, the two lowering identities, \eqref{eq:mult}, and their
ordinary differential consequences.
\end{proposition}
\begin{proof}
Fix one $\delta\in\Delta$. For
$f=L_{1,0}(x+\delta)$ and $g=L_{0,1}(x+\delta)$,
\[
 U=-e^{-\delta}\frac{f}{1+f},\qquad
 V=-e^{-\delta/\qh}\frac{\qh g}{1+\qh g}.
\]
These identities of meromorphic functions prove $F_\qh\subset K_{\qh,\Delta}$;
the denominators are not the zero functions. Each $J_{b,\rho}$ is an allowed
lattice translate, and so are its derivatives. Thus the right side of
\eqref{eq:tower-generation} is contained in the left side.
Call the right side $F_*$. It is a differential field by construction.

All axis letters belong to $F_\qh$, by \eqref{eq:axis-product} and rescaling.
Every $L_{1,b}(x+\eta)$ with $\eta\in\Delta+\Lambda_\qh$ is
$\sigma^n\tau^mJ_{b,\rho}$ for some $n,m\in\Z$ and
$\rho\in\mathscr R_\Delta$. The $b$-dependent offset in the tower is
subtracted before choosing its representative modulo $2\Lambda_\qh$.
Equations \eqref{eq:tower-sigma} and \eqref{eq:tower-tau} show that $F_*$ is
preserved by both shifts and their inverses. More explicitly,
\[
 \sigma^{-1}J_{b,\rho}
 =J_{b,\rho}-\sigma^{-1}(A_{b,\rho}R_\rho),\qquad
 \tau J_{b,\rho}=J_{b,\rho}+\tau J_{b-1,\rho}.
\]
The last equation is finite by induction on $b$, starting with
$J_{0,\rho}\in F_\qh$. Differentiating all these identities treats the jet
generators. Thus all translates of the $a=1$ mixed letters lie in $F_*$.

For $a\geq1$ and $b\geq1$, combine \eqref{eq:mult} at $z+\pi\ii$ with
\eqref{eq:lowerX} for $L_{a+1,b}$ to obtain
\begin{equation}\label{eq:reduce-a}
 \begin{split}
 L_{a+1,b}(z)
 ={}&\frac{z-\pi\ii(a-1)}{2\pi\ii a}L_{a,b}(z+\pi\ii)\\
 &-\frac{\qh b}{2a}\bigl(
 L_{a,b+1}(z+\pi\ii+\pi\ii\qh)
 +L_{a,b+1}(z+\pi\ii-\pi\ii\qh)\bigr).
 \end{split}
\end{equation}
Induction on the first index reduces every mixed letter to the preceding
case. Each step lowers that index, although the second index can increase;
hence the reduction is finite. The new offsets are still in
$\Delta+\Lambda_\qh$. Since $F_*$ is differential, all derivatives also
belong to it. This proves the reverse inclusion.
\end{proof}

\begin{proof}[Proof of \cref{thm:maincoeff}]
\Cref{prop:tower-generate} proves generation.
\Cref{lem:fixedfield,cor:tower-independent} prove that $U,V$ and all tower
derivatives are jointly algebraically independent over $\C(x)$.
Consequently evaluation identifies the indicated rational function field
with $K_{\qh,\Delta}$. The differential structure is given by
\eqref{eq:D-base} and $D(D^jJ_{b,\rho})=D^{j+1}J_{b,\rho}$.
\end{proof}

\begin{example}[One offset and the two terminal corrections]
For $\Delta=\{0\}$ one may take
\[
 \mathscr R_\Delta=\{0,\pi\ii,\pi\ii\qh,\pi\ii(1+\qh)\}.
\]
For every $b\geq1$ these give four differentially independent towers; all
towers and all their derivatives are independent jointly over $F_\qh$.
For each representative,
\[
 \sigma J_{1,\rho}-J_{1,\rho}=\qh^{-1}R_\rho(V),\qquad
 \tau J_{1,\rho}-J_{1,\rho}=L_{1,0}(x+\rho+\pi\ii\qh).
\]
Thus the towers are not themselves periodic. The elementary corrections in
both directions are essential, especially for the case $b=1$.
\end{example}

\begin{remark}[Exact scope of the nonlinear presentation]
The pole calculation is in $\C(U)(x)(V)$; it does not rely on $|\mathfrak q|$
being different from $1$. Infinite order is sufficient. Irrationality is
nevertheless indispensable. At $\qh=1$ the integral gives
$J_{1,0}=L_{1,1}=L_{2,0}\in\C(x,e^x)$, so the independence statement is false.
Nor is independence asserted after adjoining arbitrary single-valued
meromorphic functions: those could include the $J_{b,\rho}$ themselves.
The coefficient-field presentation is precisely over $F_\qh$.
\end{remark}

\section{The Laplace--simplex and Chen identification}\label{sec:laplace}

This section only requires $\qh>0$, not irrationality. Let
$a_j,b_j\geq0$, $a_j+b_j>0$, and $n_j\geq1$.

\begin{definition}[Several-variable upper prescription]
In \eqref{eq:F-intro}, use $p_j\in C_{\eta_j}$ with
$0<\eta_j<\eta_{\qh}$. The notation $(\R+\ii0)^m$ means the value obtained
with positive heights, continued within that positive-height chamber towards
zero. It is not the product of independent real-axis principal values.
\end{definition}

\begin{lemma}\label{lem:multiheight}
On the product strip
\begin{equation}\label{eq:strip}
 |\Im\omega_j|<c_{a_j,b_j}\qquad(1\leq j\leq m),
\end{equation}
\eqref{eq:F-intro} is absolutely convergent on every fixed set of positive
heights and is independent of those heights.
\end{lemma}
\begin{proof}
For $q_j=p_1+\cdots+p_j$,
$\Im q_j=\eta_1+\cdots+\eta_j>0$, so
$|q_j|^{-n_j}\leq(\eta_1+\cdots+\eta_j)^{-n_j}$.
The remaining factors have product exponential decay in the real momentum
variables by the proof of \cref{lem:height}. Vary one height at a time.
All cumulative denominators involving that variable retain positive imaginary
part throughout the deformation, and the hyperbolic-sine factors have no poles
in the intervening strip. Rectangle estimates prove independence.
\end{proof}

\begin{theorem}[Laplace--simplex identity]\label{thm:laplace}
For $\boldsymbol\omega$ in \eqref{eq:strip},
\begin{equation}\label{eq:laplace-simplex}
 \begin{split}
 \calF_{\mathbf a,\mathbf b,\mathbf n}^{\qh}(\boldsymbol\omega)
 =\int_{s_1>\cdots>s_m>0}
 \prod_{j=1}^m L_{a_j,b_j}(\omega_j-s_j)
 \frac{(s_j-s_{j+1})^{n_j-1}}{(n_j-1)!}
 \dd s_1\cdots\dd s_m,
 \qquad s_{m+1}=0.
 \end{split}
\end{equation}
The integral on the right is absolutely convergent and locally uniformly so.
Both sides determine the same continued germ wherever the left side is
continued along a common path. In particular, the right side gives a
holomorphic continuation to $\Re\omega_j<0$ for all $j$.
\end{theorem}
\begin{proof}
Take fixed positive heights and put $q_j=p_1+\cdots+p_j$.
For $n\geq1$ and $\Im q>0$,
\begin{equation}\label{eq:Laplaceone}
 q^{-n}=\frac{(-\ii)^n}{(n-1)!}
 \int_0^\infty t^{n-1}e^{\ii qt}\dd t.
\end{equation}
Insert this identity for each $q_j^{-n_j}$. On a compact substrip the absolute
value of the resulting integrand is bounded by a constant times
\begin{equation}\label{eq:Fubini-bound}
 \prod_{j=1}^m e^{-\epsilon_j|\Re p_j|}
 \prod_{j=1}^m t_j^{n_j-1}
 e^{-(\eta_1+\cdots+\eta_j)t_j},\qquad \epsilon_j>0.
\end{equation}
This is integrable in every real momentum and every $t_j>0$, so Fubini applies.
Set $s_j=t_j+\cdots+t_m$. This has Jacobian $1$ and transforms the positive
orthant into the ordered simplex in \eqref{eq:laplace-simplex}. Moreover,
\[
 \sum_jq_jt_j=\sum_jp_js_j,
 \qquad
 \ii^{|\mathbf n|-m}(-\ii)^{|\mathbf n|}=\ii^{-m}.
\]
The separate $p_j$-integrals are
$\ii L_{a_j,b_j}(\omega_j-s_j)$ by \cref{lem:height}. Their factors of $\ii$
cancel the prefactor, proving the formula for the actual positive-height
prescription.

For convergence of the right side, \cref{lem:height} on fixed $C_\eta$ gives
$|L_{a_j,b_j}(\omega_j-s_j)|\leq C e^{-\eta s_j}$ on compact substrips.
This dominates every polynomial gap factor. To reach
$\Re\omega_j<0$ with arbitrary imaginary parts, use \cref{lem:decay} for the
large $s_j$ and holomorphy of each letter in the left half-plane for bounded
$s_j$. These estimates are locally uniform, giving a holomorphic integral on
the product of left half-planes. It agrees with the original function on the
nonempty overlap. Uniqueness of analytic continuation proves the final
assertion about common paths; it does not assert that a fixed undeformed
simplex represents every continued branch.
\end{proof}

\begin{corollary}[An ascending Chen presentation]\label{cor:Chen}
Write $\omega_j=x+\delta_j$ with fixed offsets, allowing lattice shifts to be
absorbed in them. Put $\theta_j(t)=L_{a_j,b_j}(t+\delta_j)\dd t$ and $e=\dd t$.
On a sufficiently far left half-plane,
\begin{equation}\label{eq:Chen-weighted}
 \calF_{\mathbf a,\mathbf b,\mathbf n}^{\qh}
 (x+\delta_1,\ldots,x+\delta_m)
 =I[\theta_1|e^{n_1-1}|\theta_2|e^{n_2-1}|\cdots|
       \theta_m|e^{n_m-1}](x).
\end{equation}
Here $e^k$ means $k$ consecutive letters, and $I$ uses ascending order from
$-\infty$ to $x$. In particular, when all $n_j=1$ the right side is the pure
word $I[\theta_1|\cdots|\theta_m]$.
\end{corollary}
\begin{proof}
Set $t_j=x-s_j$ in \eqref{eq:laplace-simplex}. Then
$-\infty<t_1<\cdots<t_m<x$, and, with $t_{m+1}=x$,
$s_j-s_{j+1}=t_{j+1}-t_j$. Finally
\[
 \frac{(v-u)^k}{k!}
 =\int_{u<u_1<\cdots<u_k<v}\dd u_1\cdots\dd u_k.
\]
This inserts exactly the displayed blocks of $e$'s. For complex $x$ all
variables run along the same horizontal ray; the parametrization by the
real $s_j$ fixes the orientation without an additional sign.
\end{proof}

\begin{remark}
The letters $\theta_j$ in \eqref{eq:Chen-weighted} are quantum Fourier
transforms, not rational algebraic differentials. Thus this presentation is
not an expression of irrational-parameter quantum polylogarithms as ordinary
mixed-Tate iterated integrals. It gives a direct transform-level explanation
of the shuffle structure recorded in \cite[Theorem~3.7]{Goncharov2026}.
\end{remark}

\begin{proposition}[Elimination of the polynomial gap blocks]\label{prop:weight-reduction}
Every function on the left of \eqref{eq:Chen-weighted} is a finite
$\C[x]$-linear combination of pure words with letters in
$\calW_{\qh,\Delta}$. Every pure word in this space is a finite constant-linear
combination of quantum polylogarithms with all weights $n_j=1$.
\end{proposition}
\begin{proof}
Rather than integrate exact letters separately from an infinite endpoint,
expand the polynomial
\[
 \prod_{j=1}^m\frac{(t_{j+1}-t_j)^{n_j-1}}{(n_j-1)!},\qquad t_{m+1}=x,
\]
as a finite sum of monomials $c x^{k_0}\prod_jt_j^{k_j}$.
Each $t_j^{k_j}L_{a_j,b_j}(t_j+\delta_j)$ is a finite constant-linear
combination of translated letters by \cref{prop:polyclosure}. Substitution
into the absolutely convergent simplex integral gives the first claim.
For the converse use multilinearity and the case $n_1=\cdots=n_m=1$ of
\cref{cor:Chen}.
\end{proof}

\section{The shuffle differential algebra and word independence}\label{sec:bar}

We give the algebraic argument in the generality needed here, distinguishing
a chosen letter space from the whole de~Rham quotient of the coefficient
field. Let $(K,D)$ be a differential field of characteristic zero with
constant field $C$, and let $V\subset K$ be a $C$-vector subspace such that
\begin{equation}\label{eq:Vcondition}
 V\cap DK=0.
\end{equation}
Equivalently, the map $V\to K/DK$ is injective.

\begin{definition}[Shuffle differential algebra]\label{def:bar}
Set $\calB(K,V)=K\otimes_C\Sh_C(V)$. After choosing a basis
$\{f_\alpha\}_{\alpha\in\mathscr A}$ of $V$, write its word tensors as
$[\alpha_1|\cdots|\alpha_n]$. Extend $D$ by
\begin{equation}\label{eq:barD}
 D_{\calB}\bigl(g[\alpha_1|\cdots|\alpha_n]\bigr)
 =(Dg)[\alpha_1|\cdots|\alpha_n]
 +g f_{\alpha_n}[\alpha_1|\cdots|\alpha_{n-1}],
\end{equation}
with the second term omitted for the empty word.
The associated two-term differential complex is
$\calB\xrightarrow{D_{\calB}}\calB\dd x$.
\end{definition}

The formula is a derivation for shuffle: the last letter of a shuffled word
comes from exactly one of its two factors. It defines the usual
one-variable reduced-bar model for the selected alphabet. The constant word
factor $\Sh_C(V)$ has its deconcatenation Hopf structure; $\calB$ is a
commutative differential $K$-algebra and a left comodule algebra over that
Hopf algebra via
\begin{equation}\label{eq:coaction}
 g[\alpha_1|\cdots|\alpha_n]\longmapsto
 \sum_{j=0}^n[\alpha_1|\cdots|\alpha_j]\otimes
 g[\alpha_{j+1}|\cdots|\alpha_n].
\end{equation}
Here differentiation acts only in the second factor. We do not assert that
\eqref{eq:barD} is a Hopf-algebra derivation for a coproduct differentiated in
both factors. No canonical choice of representatives of all of $K/DK$ is
being made.

\begin{theorem}[Independence of iterated words]\label{thm:wordind}
Suppose an extension differential field realizes functions $I_w$ indexed by
words in $\mathscr A$, with $I_{\emptyword}=1$ and
\begin{equation}\label{eq:wordderivative}
 D I_{v\alpha}=f_\alpha I_v.
\end{equation}
Assume the coefficient field $K$ embeds in this extension and the
$f_\alpha$ represent $C$-linearly independent classes in $K/DK$.
Then the $I_w$ are linearly independent over $K$.
If they satisfy the shuffle identities, evaluation
$\calB(K,V)\to K[I_w]$ is an isomorphism of differential $K$-algebras.
\end{theorem}
\begin{proof}
Suppose a nonzero formal relation $\sum_wg_wI_w=0$ exists, with finite
support. Choose one minimizing the maximal word length $n$, then minimizing
the number of its nonzero length-$n$ coefficients. Length zero is impossible
because $K$ embeds. Divide by a chosen top coefficient so that it becomes $1$.

Differentiate the relation. By \eqref{eq:wordderivative}, its top coefficients
are $Dg_w$, and one of them is zero. If any top coefficient is nonzero, the
result is a nonzero relation with fewer top words, a contradiction. Thus all
top coefficients are constants $c_w\in C$. The differentiated relation now
has length at most $n-1$, so by minimality it is formally zero. For each
prefix $v$ of length $n-1$ its coefficient is
\[
 Dg_v+\sum_{\alpha}c_{v\alpha}f_\alpha=0.
\]
The classes of the $f_\alpha$ are independent in $K/DK$, hence all
$c_{v\alpha}$ vanish. This contradicts the choice of the top coefficient.
The first assertion follows. Shuffle makes evaluation an algebra map, and
word independence proves its injectivity; surjectivity is its definition.
\end{proof}

\begin{remark}
This minimal-length proof is the unipotent iterated-integral counterpart of
the familiar antiderivative-independence argument. It is given in full so that
no unstated differential-Galois hypothesis is required. The bar and
iterated-integral background is classical; see \cite{Chen1977,Brown2009}.
In particular the argument does not use an arithmetic theorem about values.
\end{remark}

\begin{corollary}[Constants and Lyndon generators]\label{cor:lyndon}
The constants of $(\calB(K,V),D_{\calB})$ are $C$. For any total order on a
basis of $V$, its shuffle algebra is polynomial on the Lyndon words. Under
\cref{thm:wordind}, the corresponding iterated integrals are algebraically
independent over $K$.
\end{corollary}
\begin{proof}
If $D_{\calB}F=0$, its longest word coefficients are constants. The
coefficients one level lower, followed by \eqref{eq:Vcondition}, make every
positive-length top coefficient zero. Repeating leaves $F\in C$.
The polynomial presentation of a characteristic-zero shuffle algebra is
Radford's theorem \cite{Radford1979}; its application to iterated-integral
algebras is also recalled in \cite[Section~6]{Brown2009}. Any finite
polynomial relation here involves a finite subalphabet, so the usual
finite-alphabet theorem suffices even when $V$ is infinite-dimensional.
\end{proof}

\section{Functional completeness and absolute normal forms}\label{sec:completeness}

Fix irrational $\qh>0$ and finite $\Delta$. Let
$K=K_{\qh,\Delta}$ and $V=\calW_{\qh,\Delta}$. By
\cref{thm:mainletters}, $V\cap DK=0$ and $V\cong\calV_{\qh,\Delta}$.
Because $K\subset\Mer(\C)$ contains $\C$, its differential constants are
exactly $\C$.

Choose a complex basis of $V$, for example the explicit one obtained from
\cref{prop:basis} in each offset coset. On a half-plane
\[
 \Re x<-1-\max_{\delta\in\Delta}\Re\delta,
\]
all letters are holomorphic. Any finite family decays exponentially along
leftward horizontal rays by \cref{lem:decay}. Thus all finite words
\eqref{eq:Chen-intro} converge locally uniformly and are holomorphic there.
One may verify holomorphy and path independence by first starting at $-R$
and then sending $R\to\infty$; vertical segments at $\Re x=-R$ have vanishing
contribution. Differentiating at the endpoint gives
\eqref{eq:wordderivative} with the ascending convention. The ordinary
partition of the product of two simplices gives
\begin{equation}\label{eq:shuffleidentity}
 I[u](x)I[v](x)=\sum_{w\in u\shuffle v}I[w](x),
\end{equation}
including multiplicities.

\begin{proof}[Proof of \cref{thm:mainshuffle}]
\Cref{thm:wordind} gives an injective differential algebra map
\[
 K\otimes_\C\Sh(V)\longrightarrow K[I_w].
\]
It is surjective by construction. By \cref{prop:weight-reduction},
$K[I_w]$ equals the algebra generated by all positive-integral-weight quantum
polylogarithms on the specified common-translation lines. Substituting
$V\cong\calV_{\qh,\Delta}$ proves \eqref{eq:mainshuffle}.
The statement about Lyndon words follows from \cref{cor:lyndon}.
\end{proof}

\subsection{The meaning of complete relations}

The proof supplies a finite reduction procedure for any finite proposed
relation. Transform each positive-weight generator by
\eqref{eq:laplace-simplex}; expand its polynomial gaps and use
\eqref{eq:mult} to remove polynomial multiples of letters. Combine offsets
in the same lattice coset. Reduce each letter by the partial-fraction basis
of \cref{prop:basis}, then reduce products by shuffle. The result is a
finite expression
\begin{equation}\label{eq:normalform}
 \sum_w g_w(x)I[w](x),\qquad g_w\in K,
\end{equation}
where the words are in the chosen basis of $\calV_{\qh,\Delta}$.
The original expression vanishes if and only if every $g_w$ is zero in $K$.
By \cref{prop:tower-generate,thm:maincoeff}, each coefficient reduces further
to a rational expression in $x,U,V,D^jJ_{b,\rho}$. Clearing denominators
reduces its vanishing to a polynomial identity in algebraically independent
variables. This last step, absent from a merely relative presentation, is now
justified by \cref{thm:maincoeff}.

The reductions consist of constant-linearity, elementary axis evaluation,
the two modular letter relations, Fourier integration by parts, ordinary
differentiation and rational arithmetic, the Laplace--Chen conversion, and
shuffle. They give a complete finite symbolic reduction for every finite
expression in the stated generators. This is an algebraic statement over
$\C$: the offsets, $\qh$, the representative choices, and their complex
constants are treated as exact data. It is not a numerical decision algorithm
for equality of arbitrary complex inputs or for membership in
$\Lambda_\qh$.

\begin{corollary}[Larger single-valued coefficient fields]\label{cor:largerfield}
The word-independence assertion remains valid over any differential field
$K'$ with
$K_{\qh,\Delta}\subseteq K'\subseteq\Mer(\C)$.
\end{corollary}
\begin{proof}
\Cref{thm:mainletters} excludes primitives in all of $\Mer(\C)$, and the
constant field of $K'$ is still $\C$. Apply \cref{thm:wordind}.
\end{proof}

\subsection{Combining coefficient and shuffle independence}
\label{sec:absolute-completeness}

\begin{proof}[Proof of \cref{cor:mainabsolute}]
The coefficient evaluation in \cref{thm:maincoeff} is an isomorphism onto
$K_{\qh,\Delta}$. By \cref{thm:mainshuffle,cor:lyndon}, adjoining the Lyndon integrals is a
polynomial extension over that field. Composing these two isomorphisms gives
\eqref{eq:mainabsolute}. In particular, a polynomial relation among the
combined generators can be regarded as a polynomial in the Lyndon variables
with coefficients in $K_{\qh,\Delta}$. Their independence makes every such
coefficient zero; coefficient-field independence then makes every original
coefficient polynomial zero. This proves joint independence.
\end{proof}

The differential structure in this explicit presentation is also determined:
\[
 Dx=1,\qquad DU=U,\qquad DV=V/\qh,\qquad
 D(D^jJ_{b,\rho})=D^{j+1}J_{b,\rho},
\]
and differentiation of words is \eqref{eq:wordderivative}. The letter
coefficient on its right has an explicit tower expression by
\cref{prop:tower-generate}; shuffle reduction then expresses the result in
Lyndon generators. Hence derivatives of positive-weight generators introduce
no unspecified field of coefficients either.

\begin{corollary}[Positive weights without adjoining the coefficient field]
\label{cor:positive-absolute}
Let $A_{\qh,\Delta}^{>0}$ be the $\C(x)$-algebra generated by the same
positive-weight quantum polylogarithms, without adjoining $K_{\qh,\Delta}$.
Then
\[
 A_{\qh,\Delta}^{>0}\cong
 \C(x)\otimes_\C\Sh(\calV_{\qh,\Delta}).
\]
\end{corollary}
\begin{proof}
Restrict the injective map in \cref{thm:mainshuffle} to coefficients in
$\C(x)$. Its image is the stated algebra by
\cref{prop:weight-reduction}. In particular surjectivity uses only polynomial
coefficients in $x$, not arbitrary elements of $K_{\qh,\Delta}$.
\end{proof}

\begin{remark}
\Cref{cor:positive-absolute} already follows from the relative shuffle theorem.
The stronger content of \cref{cor:mainabsolute} is that nonlinear identities
inside the adjoined weight-zero differential field have now been classified,
not that the positive-weight coefficient field was simply made smaller.
\end{remark}

\subsection{Small examples and normalization checks}

\begin{example}[The classical one-frequency sector]
The first letters are
\begin{equation}\label{eq:classicalexamples}
 \begin{split}
 L_{1,0}(x)&=-\frac{e^x}{1+e^x},\\
 L_{2,0}(x)&=\frac{x}{2\pi\ii}\frac{e^x}{1-e^x},\\
 L_{3,0}(x)&=\frac{x^2+\pi^2}{2(2\pi\ii)^2}
               \frac{-e^x}{1+e^x}.
 \end{split}
\end{equation}
They follow from \eqref{eq:axesbase} and \eqref{eq:axisrecursion} and agree
with \cite[Lemma~5.4]{Goncharov2026}. In particular,
\[
 I[L_{1,0}](x)=-\log(1+e^x),\qquad
 I[L_{1,0}|L_{1,0}](x)=\frac12\log^2(1+e^x),
\]
with the branches fixed at $-\infty$. The factor $1/2$ is the shuffle
multiplicity, not an extra relation in the letter space.
\end{example}

\begin{example}[Eliminating one additional weight]
For $\omega=x+\delta$, \cref{thm:laplace} and \eqref{eq:mult} give
\begin{equation}\label{eq:weighttwo}
 \begin{split}
 \calF_{a,b,2}^{\qh}(\omega)
 ={}&\omega\calF_{a,b,1}^{\qh}(\omega)\\
 &-\ii\pi a\left(
 \calF_{a+1,b,1}^{\qh}(\omega+\ii\pi)
 +\calF_{a+1,b,1}^{\qh}(\omega-\ii\pi)\right)\\
 &-\ii\pi\qh b\left(
 \calF_{a,b+1,1}^{\qh}(\omega+\ii\pi\qh)
 +\calF_{a,b+1,1}^{\qh}(\omega-\ii\pi\qh)\right).
 \end{split}
\end{equation}
The constants are fixed by decay at $-\infty$. This identity illustrates
why the depth of a pure-word presentation need not equal Goncharov's
weight: the kernel orders $a,b$ can increase when a polynomial gap is
removed.
\end{example}

\begin{example}[The two-frequency letter]
For irrational $\qh$, $L_{1,1}(z)\dd z$ has residue $-1$ at every
$z=\ii\pi(1+2m+\qh(1+2n))$ and residue $+1$ at the corresponding negative
point. Its reflection identity is
\[
 L_{1,1}(z)-L_{1,1}(-z)=\frac{\ii z}{2\pi\qh}.
\]
The nonzero residues exclude a single-valued meromorphic primitive, even if
one enlarges the coefficient field far beyond $K_{\qh,\Delta}$.
\end{example}

\section{Scope, functoriality, and remaining problems}\label{sec:limits}

\subsection{The letter quotient is not the whole differential cohomology}

The injection
\[
 \calV_{\qh,\Delta}\hookrightarrow K_{\qh,\Delta}\dd x/dK_{\qh,\Delta}
\]
is generally not surjective. Choose $z_0$ outside the finite union
$\bigcup_{\delta\in\Delta}(-\delta+\Lambda_{\qh})$.
Then $\dd x/(x-z_0)$ belongs to $K_{\qh,\Delta}\dd x$ and has a nonzero
residue at $z_0$, while every quantum letter is regular there. It cannot be
cohomologous to any finite quantum-letter combination. Accordingly
$\calB(K,V)$ here is the shuffle extension generated by a selected
cohomologically independent alphabet, not a universal primitive closure of
$K$. We make no assertion that its first de~Rham cohomology vanishes.

\subsection{Functional variables and continuation}

Both completeness statements fix all offsets and vary one common translation
$x$. They apply to every such fixed offset configuration, after lattice
cosets are identified. It does not supply a universal algebraic presentation
with the offsets themselves as varying coefficient parameters. Similarly,
meromorphic continuation of the identities preserves equality of continued
germs, but need not preserve independence after evaluating $x$ at a number.
A numerical or arithmetic specialization theorem would require new input.

The nonlinear theorem concerns the differential field generated by the
\emph{depth-one} weight-zero letters; the iterated-integral theorem concerns
positive integral weights at arbitrary depth. These statements do not claim
an independent classification for arbitrary mixed positive, zero, and
negative weight indices in the original multiple-momentum definition.

The Laplace--Chen identity holds for every positive real $\qh$. The explicit
residue quotient and both independence arguments use irrationality. When
$\qh=r/s\in\Q$, pole labels collide; the injectivity of
$(k,\ell)\mapsto k+\qh\ell$ fails. The rational-parameter theory should not
be inferred by substituting into the irrational-parameter independence
statement. For example, at $\qh=1$ the kernels themselves give
$L_{1,1}=L_{2,0}$, whereas
$d_X^{-1}d_Y^{-1}-d_X^{-2}$ is nonzero in $\calS/\calR$.
Thus neither the irrational letter presentation nor the differential tower
independence can remain unchanged at rational parameter. Goncharov treats reduction to multiple polylogarithms at rational
parameter in \cite[Theorems~3.2--3.3]{Goncharov2026}.

\subsection{Functoriality requires actual differential data}

For a differential-field isomorphism $\phi:K\to K'$ carrying a specified
letter space $V$ onto $V'$ and commuting with derivations, there is a
letterwise differential-algebra isomorphism
\[
 \calB(K,V)\longrightarrow\calB(K',V'),\qquad
 g[f_1|\cdots|f_n]\longmapsto
 \phi(g)[\phi(f_1)|\cdots|\phi(f_n)].
\]
This statement concerns the \emph{actual representative forms}. If only
cohomology classes are transported, representative changes by exact forms
can produce lower-length integration-by-parts and endpoint terms. The
associated graded action alone does not determine the normalized action on
iterated integrals.

Moreover a quantum-cluster transformation acts on a noncommutative quantum
torus, not automatically on the commutative differential field
$K_{\qh,\Delta}$. The operator theorem below proves a faithful presentation
of a specified operator algebra. It does not assert that conjugation by the
pentagon intertwiner carries scalar multiplication by every Fourier letter
to scalar multiplication by another such letter. Consequently no universal
cross-chart scalar completeness theorem is deduced here from the operator
presentation.

\subsection{Arithmetic motivation}

Radchenko--Wheeler construct finite quantum dilogarithms from special values
of a modular extension of the quantum dilogarithm and prove algebraicity
\cite{RadchenkoWheeler2026}. Their second version gives additional algebraicity
arguments and explicitly distinguishes the small-order normalization exceptions.
Kopp interprets real-quadratic Stark class invariants using the
Shintani--Faddeev modular cocycle \cite{Kopp2024}.
Appleby--Flammia--Kopp's earlier work gives the conditional Stark/ghost-SIC
and rank-$r$ framework \cite{AFK2025}; their later paper proves the
all-admissible-rank twisted convolution identity and supplies the
RW/Shintani--Faddeev dictionary \cite[Theorem~1.3 and
Proposition~2.1]{ApplebyFlammiaKopp2026}. These finite identities and their
arithmetic applications are background, not conclusions of this paper.

The shared analytic object motivates asking how a functional relation theory
interacts with these special values. A literal description as fibers would
require a defined family or specialization functor and a proof of its kernel.
Neither the scalar theorem above nor the operator theorem below supplies
that arithmetic comparison. Our companion note \cite{LongFinite2026}
studies the raw finite equations, with explicit low-order certificates and
effective degree bounds, and finite \'{e}taleness of the raw solution scheme
through a dual-number realization and invariant scalar reconstruction. None of those assertions is an input
to the present functional or operator proofs.

\appendix
\section{Faithfulness of the pentagon crossed product}\label{app:operator}

This appendix proves the operator statement independently of the scalar
completeness theorem. Its analytic inputs are Goncharov's invariant Schwartz
space and covariance and pentagon theorems \cite[Theorems~1.1, 2.3,
2.6]{Goncharov2007}. No inverse difference operator will be applied to a
vector.

\subsection{The Ore-denominator lemma}

\begin{theorem}[Independence of intertwiner powers]\label{thm:oreind}
Let $A$ be a left Ore domain over $\C$, let $D=\Frac(A)$ be its fraction
division ring, and let $\sigma\in\Aut(A)$ extend to $D$. Suppose
$\sigma^j$ is outer on $D$ for $1\leq j<m$.
Let $\pi:A\hookrightarrow\End_\C(V)$ be faithful, and suppose
$U\in\operatorname{GL}_\C(V)$ satisfies
\begin{equation}\label{eq:covgeneral}
 U\pi(a)=\pi(\sigma(a))U\qquad(a\in A).
\end{equation}
No action of $D$ on $V$ is required. Then
\begin{equation}\label{eq:oreind}
 \sum_{j=0}^{m-1}\pi(a_j)U^j=0\quad\Longrightarrow\quad a_j=0\text{ for all }j.
\end{equation}
\end{theorem}
\begin{proof}
Choose a nontrivial relation $R=\sum_j\pi(a_j)U^j=0$ with the fewest nonzero
coefficients. Right multiplication by the inverse of the smallest occurring
power of $U$ arranges $a_0\ne0$, with exponents still between $0$ and $m-1$.
There are at least two nonzero coefficients, by faithfulness and invertibility
of $U$.

For any $a\in A$, choose $s\in A\setminus\{0\}$ and $t\in A$ such that
\[
 s a_0aa_0^{-1}=t\quad\hbox{in }D,
 \qquad\hbox{hence }s a_0a=t a_0.
\]
This is permitted by the left-fraction description of $D$. The relation
$\pi(s)R\pi(a)-\pi(t)R=0$ has coefficient zero at $U^0$ and reads
\[
 \sum_j\pi\bigl(s a_j\sigma^j(a)-t a_j\bigr)U^j=0.
\]
If any of its coefficients is nonzero \emph{in $A$}, it contradicts minimal
support. Therefore $s a_j\sigma^j(a)=t a_j$ in $A$ for every $j$.
Only now, in the abstract division ring, cancel $s$ and write
$d_j=a_0^{-1}a_j$. The result is
\[
 d_j\sigma^j(a)=a d_j\qquad(a\in A).
\]
For an active $j>0$, $d_j\ne0$, and this makes
$\sigma^j(a)=d_j^{-1}ad_j$ on $A$, hence on its fraction division ring.
This contradicts outerness. No active $j>0$ remains, contrary to the choice
of a relation with at least two terms.
\end{proof}

\begin{remark}
The argument never cancels the operator $\pi(s)$, which could have a kernel.
It first obtains an equality of coefficients in $A$, and only then uses
inverses in $D$. This distinction is the reason analytic Ore localization is
unnecessary.
\end{remark}

\subsection{The modular-double coefficient algebra}

Fix real irrational $\qh>0$ and set $q=e^{\pi\ii\qh}$,
$q^\vee=e^{\pi\ii/\qh}$. All coefficient algebras in this appendix are
specialized over $\C$ at these values. Let $T$ be the tensor product of the
two quantum tori with generators
$X^{\pm1},Y^{\pm1},(X^\vee)^{\pm1},(Y^\vee)^{\pm1}$ and relations
\begin{equation}\label{eq:qtori}
 XY=q^2YX,\qquad X^\vee Y^\vee=(q^\vee)^2Y^\vee X^\vee,
\end{equation}
with the two pairs commuting. Ordered Laurent monomials are a basis.

Let $\calL$ be the specialized coefficient algebra denoted $L$ in
\cite[Section~2.2]{Goncharov2007}. We use two explicit properties established
there: $\calL$ is preserved by the order-five automorphism $\gamma$ and
\begin{equation}\label{eq:P-in-L}
 P=\C\langle X^{-1},Y,(X^\vee)^{-1},Y^\vee\rangle
 \subseteq\calL\subseteq T.
\end{equation}
The first canonical-basis cone consists, up to nonzero powers of $q$, of
$X^aY^b$ with $a\leq0$, $b\geq0$, and similarly in the dual factor;
this proves the first inclusion. On the fraction division ring, the first
factor of $\gamma$ is
\begin{equation}\label{eq:gamma}
 \gamma(X)=Y^{-1},\qquad \gamma(Y)=(1+qY)X,
\end{equation}
and the dual factor has the analogous formula with $q^\vee$.
Goncharov proves $\gamma^5=1$.

\begin{lemma}\label{lem:OreP}
$P$, $\calL$, and $T$ are left and right Ore domains with the same fraction
division ring $D$.
\end{lemma}
\begin{proof}
Products of ordered monomials differ from the corresponding ordered
monomial only by a nonzero scalar; leading-term comparison shows that all
three algebras are domains. Let $P_{\leq n}$ be the span of monomials in its
four generators of total degree at most $n$. Then
$\dim P_{\leq n}=\binom{n+4}{4}$.
For nonzero $a,b\in P$ of degrees at most $d$, the subspaces
$P_{\leq n}a$ and $P_{\leq n}b$ both have this dimension and lie in
$P_{\leq n+d}$. For large $n$,
\[
 2\binom{n+4}{4}>\binom{n+d+4}{4}.
\]
Their intersection is nonzero, giving $sa=tb$ with nonzero $s,t\in P$.
The right-sided argument is identical, so $P$ is Ore.

Its fraction division ring contains the inverses of the four generators,
hence contains $T$. For $a,b\in\calL$, $a,b\ne0$, write
$ab^{-1}=s^{-1}t$ with $s,t\in P\setminus\{0\}$. Then $sa=tb$ is a common
left multiple in $\calL$; right fractions give the right Ore condition.
The same reasoning applies to $T$, and all fraction division rings equal
$\Frac(P)$.
\end{proof}

\begin{lemma}[Outerness]\label{lem:outer}
Every nonidentity power $\gamma^j$, $1\leq j\leq4$, is outer on $D$.
\end{lemma}
\begin{proof}
For a nonzero Laurent polynomial in $T$, let $\deg_X$ be its largest
$X$-exponent. Regarding $T$ as a skew Laurent-polynomial ring in $X$ over
a domain shows $\deg_X(AB)=\deg_X(A)+\deg_X(B)$. It extends to
$D^\times$ by
\[
 \deg_X(A^{-1}B)=\deg_X(B)-\deg_X(A).
\]
To check well-definedness take common left Ore denominators; their degrees
cancel by additivity. Every inner automorphism preserves this degree.
But $\deg_X(X)=1$ and $\deg_X(\gamma(X))=\deg_X(Y^{-1})=0$, so $\gamma$
is outer. If $\gamma^j$ were inner for $1\leq j\leq4$, choose $r>0$ with
$jr\equiv1\pmod5$; then $\gamma=(\gamma^j)^r$ would be inner. This is
impossible.
\end{proof}

\subsection{The analytic representation}

On Gaussian-polynomial test functions, the torus generators act by
\begin{equation}\label{eq:operators}
 \begin{aligned}
 \widehat X f(x)&=f(x+2\pi\ii\qh),&
 \widehat Y f(x)&=e^x f(x),\\
 \widehat X^\vee f(x)&=f(x+2\pi\ii),&
 \widehat Y^\vee f(x)&=e^{x/\qh}f(x).
 \end{aligned}
\end{equation}
Goncharov constructs a common invariant space $S_{\calL}$ containing these
test functions, and an invertible quantum-dilogarithm intertwiner $K$ on it,
such that
\begin{equation}\label{eq:analyticcov}
 K^{-1}\widehat A K=\widehat{\gamma(A)}\quad(A\in\calL),
 \qquad K^5=\kappa I\quad(\kappa\in\C^\times).
\end{equation}
Here a fixed nonzero scalar normalization of his intertwiner is allowed;
$\kappa$ is determined by it. We import the common-domain and covariance
assertions from \cite[Theorems~1.1, 2.3, 2.6]{Goncharov2007}.

\begin{lemma}[Faithful coefficients]\label{lem:coefffaithful}
The representation of $\calL$ on $S_{\calL}$ is faithful.
\end{lemma}
\begin{proof}
Apply a Laurent monomial
$Y^b(Y^\vee)^dX^a(X^\vee)^c$ to $f_0(x)=e^{-x^2/2}$.
Its value is
\begin{equation}\label{eq:gaussian}
 e^{2\pi^2(a\qh+c)^2}e^{-x^2/2}
 e^{(b+d/\qh-2\pi\ii(a\qh+c))x}.
\end{equation}
For irrational real $\qh$, the displayed exponents are distinct for distinct
integer quadruples $(a,b,c,d)$: compare imaginary parts first, then real
parts. Finite families of exponentials with distinct exponents are linearly
independent, as follows for example from a Vandermonde determinant of their
derivatives at zero. Thus the torus action on the Gaussian test space is
faithful, and so is the restriction to $\calL$ on $S_{\calL}$.
\end{proof}

\begin{theorem}[Pentagon crossed-product presentation]\label{thm:crossed}
For irrational real $\qh>0$, the finite algebraic operator relations in the
algebra generated by $\widehat{\calL}$ and $K,K^{-1}$ are exactly the
coefficient relations, covariance, and $K^5=\kappa I$. More precisely,
\begin{equation}\label{eq:crossed}
 \calL[u;\gamma^{-1}]/(u^5-\kappa)
 \xrightarrow{\sim}
 \operatorname{Alg}_{\C}(\widehat{\calL},K,K^{-1}),
 \qquad u\longmapsto K.
\end{equation}
In particular $\sum_{j=0}^4\widehat A_jK^j=0$ with $A_j\in\calL$ implies
$A_0=\cdots=A_4=0$.
\end{theorem}
\begin{proof}
Covariance in \eqref{eq:analyticcov} is
$K\widehat A=\widehat{\gamma^{-1}(A)}K$.
Apply \cref{thm:oreind} with $A=\calL$, $U=K$, $m=5$, and
$\sigma=\gamma^{-1}$, using \cref{lem:OreP,lem:outer,lem:coefffaithful}.
This proves independence of the five powers.
Since $\sigma^5=1$ and $\kappa$ is a nonzero scalar, $u^5-\kappa$ is central
in the skew-polynomial ring. Division by this monic polynomial gives the
unique left-coefficient normal form $\sum_{j=0}^4 A_ju^j$.
The independence just proved makes its analytic realization injective;
surjectivity follows from the definition of the generated algebra and
$u^{-1}=\kappa^{-1}u^4$.
\end{proof}

\begin{remark}[Operator scope]
\Cref{thm:crossed} concerns finite sums and products in the specified generated
algebra. It does not classify identities in an algebra enlarged by the bare
Fourier transform, arbitrary meromorphic multiplication operators, analytic
closures, or parameter-changing intertwiners. In particular it does not
identify the operator algebra with the commutative shuffle algebra of
\cref{thm:mainshuffle}.
\end{remark}

\section*{Research and computational disclosure}
The manuscript was developed with AI-assisted literature retrieval, proof
exploration, drafting, and internal checking. The author is responsible for
its mathematical claims and references. Symbolic low-order normalization
checks and numerical contour comparisons accompany the draft as supplementary
code; they are not used as proofs. This version is a research draft and has
not undergone independent external refereeing.

\begin{thebibliography}{99}

\bibitem{Brown2009}
F.~C.~S.~Brown,
\emph{Multiple zeta values and periods of moduli spaces
$\overline{\mathfrak M}_{0,n}$},
Ann. Sci. \'Ec. Norm. Sup\'er. (4) \textbf{42} (2009), no.~3, 371--489.
\href{https://www.numdam.org/item/ASENS_2009_4_42_3_371_0/}{Journal text}.

\bibitem{Chen1977}
K.-T.~Chen,
\emph{Iterated path integrals},
Bull. Amer. Math. Soc. \textbf{83} (1977), no.~5, 831--879.
\href{https://doi.org/10.1090/S0002-9904-1977-14320-6}{doi:10.1090/S0002-9904-1977-14320-6}.

\bibitem{Goncharov2007}
A.~B.~Goncharov,
\emph{Pentagon relation for the quantum dilogarithm and quantized
$M^{\mathrm{cyc}}_{0,5}$},
in \emph{Geometry and Dynamics of Groups and Spaces},
Progress in Mathematics \textbf{265}, Birkh\"auser, 2008, 415--428.
\href{https://arxiv.org/abs/0706.4054v2}{arXiv:0706.4054v2}.
Theorem and equation numbers in the text refer to this arXiv version.

\bibitem{Goncharov2026}
A.~B.~Goncharov,
\emph{Quantum polylogarithms},
preprint, 2026.
\href{https://arxiv.org/abs/2601.00472v1}{arXiv:2601.00472v1}.

\bibitem{HardouinSinger2008}
C.~Hardouin and M.~F.~Singer,
\emph{Differential Galois theory of linear difference equations},
Math. Ann. \textbf{342} (2008), no.~2, 333--377.
\href{https://doi.org/10.1007/s00208-008-0238-z}{doi:10.1007/s00208-008-0238-z}.
Erratum, Math. Ann. \textbf{350} (2011), 243--244,
\href{https://doi.org/10.1007/s00208-010-0551-1}{doi:10.1007/s00208-010-0551-1}.

\bibitem{Schneider2010}
C.~Schneider,
\emph{Parameterized telescoping proves algebraic independence of sums},
Ann. Comb. \textbf{14} (2010), no.~4, 533--552.
\href{https://doi.org/10.1007/s00026-011-0076-7}{doi:10.1007/s00026-011-0076-7}.

\bibitem{RadchenkoWheeler2026}
D.~Radchenko and C.~Wheeler,
\emph{Real quadratic fields and finite quantum dilogarithms I},
preprint, 2026, \href{https://arxiv.org/abs/2609.21892v2}{arXiv:2609.21892v2},
September 29, 2026. Version 1 appeared September 18, 2026.

\bibitem{ApplebyFlammiaKopp2026}
M.~Appleby, S.~T.~Flammia, and G.~S.~Kopp,
\emph{The twisted convolution identity and ghost $r$-SICs from finite quantum dilogarithms},
preprint, 2026, \href{https://arxiv.org/abs/2609.39192v1}{arXiv:2609.39192v1},
September 30, 2026.


\bibitem{Huang2026}
T.-C.~Huang,
\emph{Cyclic Haagerup--Izumi fusion categories at every odd order},
\href{https://arxiv.org/abs/2609.15986v1}{arXiv:2609.15986v1}, September 14, 2026.

\bibitem{GSY2026}
T.~Gannon, A.~Schopieray, and H.~Yadav,
\emph{On Haagerup--Izumi fusion categories},
\href{https://arxiv.org/abs/2609.25185v1}{arXiv:2609.25185v1}, September 21, 2026.

\bibitem{AFK2025}
M.~Appleby, S.~T.~Flammia, and G.~S.~Kopp,
\emph{A Constructive Approach to Zauner's Conjecture via the Stark Conjectures},
\href{https://arxiv.org/abs/2501.03970v2}{arXiv:2501.03970v2}, March 17, 2025.
First version: January 7, 2025.

\bibitem{Kopp2024}
G.~S.~Kopp,
\emph{The Shintani--Faddeev modular cocycle: Stark units from $q$-Pochhammer ratios},
\href{https://arxiv.org/abs/2411.06763v3}{arXiv:2411.06763v3}, May 3, 2025.
First version: November 11, 2024.

\bibitem{Seki2026}
S.-i.~Seki,
\emph{A proof of Hirose's duality conjecture},
\href{https://arxiv.org/abs/2609.40213v1}{arXiv:2609.40213v1}, September 30, 2026.

\bibitem{BGMS2026}
J.~Bl\"umlein, A.~M.~Gavrilik, O.~Mykhailiv, and C.~Schneider,
\emph{The $q$-extension of iterated integrals and nested sums in quantum field theory},
\href{https://arxiv.org/abs/2608.02702v1}{arXiv:2608.02702v1}, August 3, 2026.

\bibitem{LongFinite2026}
C.~D.~Long,
\emph{Finite Quantum Dilogarithms: Normalization Exceptions, Infinitesimal
Rigidity, and Effective Degree Bounds}, research draft, October 1, 2026.
Companion manuscript in
\href{https://github.com/long-mathematics/quantum-polylogarithms-finite-dilogarithms}{the project repository}.

\bibitem{Radford1979}
D.~E.~Radford,
\emph{A natural ring basis for the shuffle algebra and an application to group
schemes},
J. Algebra \textbf{58} (1979), no.~2, 432--454.
\href{https://doi.org/10.1016/0021-8693(79)90171-6}{doi:10.1016/0021-8693(79)90171-6}.

\end{thebibliography}
\end{document}
